%  LaTeX support: latex@mdpi.com 
%  For support, please attach all files needed for compiling as well as the log file, and specify your operating system, LaTeX version, and LaTeX editor.

%=================================================================
\documentclass[aerospace,article,submit,moreauthors]{Definitions/mdpi} 
%\documentclass[journal,article,submit,moreauthors,unicode]{Definitions/mdpi} % For some letters Palatino lacks
%\documentclass[preprints,article,submit,moreauthors]{Definitions/mdpi} 
%\documentclass[preprints,article,submit,moreauthors,unicode]{Definitions/mdpi} 
% For posting an early version of this manuscript as a preprint, you may use "preprints" as the journal. Changing "submit" to "accept" before posting will remove line numbers.
\newcommand{\mmin}{m_{\min}}

%--------------------
% Class Options:
%--------------------
%----------
% journal
%----------
% Choose between the following MDPI journals:
% aam, aamr, aapm, acbr, accountaudit, acn, acoustics, actuators, addictions, addictprev, adhesives, adolescents, aee, aep, aerobiology, aeronautics, aerospace, agriculture, agriengineering, agrochemicals, agronomy, ai, aiaerosp, aiagents, aibiol, aichem, aieng, aimater, aimed, aipa, air, aisens, aisoc, aiua, algorithms, allergies, alloys, amh, analog, analytica, analytics, anatomia, anesthres, animals, antibiotics, antibodies, antioxidants, applbiosci, appliedchem, appliedmath, appliedphys, applmech, applmicrobiol, applnano, applsci, aps, aquacj, archaeolstud, architecture, arm, arthropoda, asc, asd, asi, astronautics, astronomy, atmosphere, atoms, audiolres, automation, axioms, bacteria, batteries, bcrc, bdcc, beverages, bioactives, biochem, bioengineering, biologics, biology, biomass, biomechanics, biomed, biomedicines, biomedinformatics, biomimetics, biomolecules, biophysica, bioresourbioprod, biosensors, biosphere, biotech, birds, blockchains, bloods, blsf, brainsci, breath, bse, buildings, cancers, carbon, cardio, cardiogenetics, cardiovascmed, catalysts, cdr, cells, ceramics, challenges, chemeco, chemengineering, chemistry, chemosensors, chemproc, children, chips, cimb, civileng, cleantechnol, climate, clinbioenerg, clinpract, clockssleep, cmd, cmtr, coasts, coatings, colloids, colorants, complexities, complications, compounds, computation, computers, condensedmatter, conservation, constrmater, cosmetics, covid, crops, cryo, cryptography, crystals, csmf, ctn, culture, curroncol, cyber, dairy, data, ddc, dentistry, dermato, dermatopathology, designs, devices, dgg, dhi, diabetology, diagnostics, dietetics, digital, disabilities, diseases, diversity, dna, drones, drugreact, drylandsci, dynamics, earth, ebj, ece, ecm, eco-environsci, ecologies, edm, eesp, electricity, electrochem, electronicmat, electronics, embodiedintell, encyclopedia, endocrines, energies, eng, engproc, entomology, entropy, environments, environremediat, enzymology, epidemiologia, epigenomes, epilepsy, esa, est, fermentation, fibers, fintech, fire, fishes, fluids, foodeng, foods, forecasting, forensicsci, forests, fossstud, foundations, fractalfract, freshwater, fuels, future, futureinternet, futurepharmacol, futurephys, futuretransp, galaxies, gases, gastroent, gastrointestdisord, gastronomy, gels, genes, geochem, geographies, geohazards, geomatics, geometry, geosciences, geotechnics, geriatrics, germs, glacies, grainsci, grasses, green, greenhealth, gucdd, gynecolcancer, hardware, hazardousmatters, healthcare, healthcmater, hearts, hemato, hematolrep, hep, heritage, higheredu, highthroughput, hls, horticulturae, hospitals, humanoids, hydrobiology, hydrogen, hydrology, hydrometeorology, hydropower, hygiene, idhm, idr, iic, ijem, ijerph, ijgi, ijmd, ijms, ijns, ijom, ijp, ijpb, ijt, ijtm, ijtpp, ijtst, ime, immuno, industries, inflammj, informatics, information, infrastructures, inorganics, insects, instruments, inventions, iot, iron, j, jaestheticmed, jal, japma, jcdd, jcm, jcp, jcrm, jcs, jcto, jdad, jdb, jdcp, jdream, jemr, jeta, jfb, jfmk, jfth, jgbg, jgg, jimaging, jissn, jlpea, jmahp, jmmp, jmms, jmp, jmse, jne, jnt, jof, joitmc, joma, jop, joptm, jor, jox, jpbi, jphytomed, jpm, jsan, jtaer, jvd, jzbg, kidneydial, kinasesphosphatases, knowledge, labmed, laboratories, lam, land, langtechnol, lasers, legumes, life, lights, limnolrev, lipidology, liquids, livers, logics, logistics, lubricants, lymphatics, machines, macromol, magnetism, magnetochemistry, make, marinedrugs, maritime, materials, materproc, mathematics, mca, mcb, measurements, meb, medicina, medicines, medsci, membranes, merits, metabolites, metals, metamaterialsadv, meteorology, methane, metrics, metrology, micro, microarrays, microbiolres, microchem, microelectronics, micromachines, microorganisms, microplastics, microwave, minerals, mining, mmphys, modelling, molbank, molecules, mollusca, mps, msf, mti, multimedia, muscles, musculoskeletsci, mutation, nanoenergyadv, nanomanufacturing, nanomaterials, ncrna, ndt, network, neuroglia, neuroimaging, neurolint, neurosci, nitrogen, numbertheory, nursrep, nutraceuticals, nutrients, obesities, occuphealth, oceans, ohbm, onco, optics, optoelectronics, oral, organics, organoids, osteology, oxygen, pandemics, parasites, parasitologia, particles, pathogens, pathophysiology, pediatrrep, pets, pharmaceuticals, pharmaceutics, pharmacoepidemiology, pharmacy, phc, philosophies, photochem, photonics, photovoltaics, phycology, physchem, physics, physiologia, pjm, plants, plasma, platforms, pnr, pollutants, polymers, polysaccharides, populations, poultry, powders, ppi, precision, precisoncol, prevhealthc, proceedings, processes, prosthesis, proteomes, psf, psychiatryint, psychoactives, purification, quantumrep, quaternary, qubs, radiation, rareearths, rdt, reactions, realestate, receptors, recycling, regeneration, remotesensing, reports, reprodmed, resources, rheumato, rjpm, robotics, rsee, ruminants, safety, sanpp, sca, sci, scipharm, sclerosis, seeds, sensors, senssyst, separations, sexes, sgf, shi, signals, sinusitis, siuj, skins, smartcities, smartfish, smartgrids, sna, societies, software, soilsystems, solar, solids, spectroscj, sports, standards, stats, std, sti, stomatology, stratsediment, stresses, strokerep, superintelligence, surfaces, surgeries, susagric, suschem, susenvmgmt, sustainability, sustainfoods, sustainpackag, symmetry, synbio, systems, tae, targets, taxonomy, technologies, telecom, test, textiles, thalassrep, therapeutics, thermo, timespace, tomography, totalenviron, toxics, toxins, tph, transplantology, transportation, traumacare, traumas, treeres, tri, tropicalmed, universe, urbansci, uro, vaccines, vehicles, veinslymphat, venereology, vetsci, vibration, virtualworlds, viruses, vision, vsd, waste, water, welding, wem, wetlres, wevj, wild, wind, winestud, women, world, xscience, zoonoticdis, zoos

%---------
% article
%---------
% The default type of manuscript is "article", but can be replaced by: 
% abstract, addendum, article, benchmark, book, bookreview, briefcommunication, briefreport, casereport, changes, clinicopathologicalchallenge, comment, commentary, communication, conceptpaper, conferenceproceedings, correction, conferencereport, creative, datadescriptor, discussion, entry, expressionofconcern, extendedabstract, editorial, essay, erratum, fieldguide, hypothesis, interestingimages, letter, meetingreport, monograph, newbookreceived, obituary, opinion, proceedingpaper, projectreport, reply, retraction, review, perspective, protocol, shortnote, studyprotocol, supfile, systematicreview, technicalnote, viewpoint, guidelines, registeredreport, tutorial, giantsinurology, urologyaroundtheworld, patentsummary

% supfile = supplementary materials

%----------
% submit
%----------
% The class option "submit" will be changed to "accept" by the Editorial Office when the paper is accepted. This will only make changes to the frontpage (e.g., the logo of the journal will get visible), the headings, and the copyright information. Also, line numbering will be removed. Journal info and pagination for accepted papers will also be assigned by the Editorial Office.

%------------------
% moreauthors
%------------------
% If there is only one author the class option oneauthor should be used. Otherwise use the class option moreauthors.

%=================================================================
% MDPI internal commands - do not modify
\firstpage{1} 
\makeatletter 
\setcounter{page}{\@firstpage} 
\makeatother
\pubvolume{1}
\issuenum{1}
\articlenumber{0}
\pubyear{2026}
\copyrightyear{2026}
%\externaleditor{Firstname Lastname} % More than 1 editor, please add `` and '' before the last editor name
\datereceived{ } 
\daterevised{ } % Comment out if no revised date
\dateaccepted{ } 
\datepublished{ } 
%\datecorrected{} % For corrected papers: "Corrected: XXX" date in the original paper.
%\dateretracted{} % For retracted papers: "Retracted: XXX" date in the original paper.
%\doinum{} % Used for some special journals, like molbank
%\pdfoutput=1 % Uncommented for upload to arXiv.org
%\CorrStatement{yes}  % For updates
%\longauthorlist{yes} % For many authors that exceed the left citation part
%\IsAssociation{yes} % For association journals

%=================================================================
% Add packages and commands here. The following packages are loaded in our class file: fontenc, inputenc, calc, indentfirst, fancyhdr, graphicx, epstopdf, lastpage, ifthen, float, amsmath, amssymb, lineno, setspace, enumitem, mathpazo, booktabs, titlesec, etoolbox, tabto, xcolor, colortbl, soul, multirow, microtype, tikz, totcount, changepage, attrib, upgreek, array, tabularx, pbox, ragged2e, tocloft, marginnote, marginfix, enotez, amsthm, natbib, hyperref, cleveref, scrextend, url, geometry, newfloat, caption, draftwatermark, seqsplit
% cleveref: load \crefname definitions after \begin{document}

%=================================================================
% Please use the following mathematics environments: Theorem, Lemma, Corollary, Proposition, Characterization, Property, Problem, Example, ExamplesandDefinitions, Hypothesis, Remark, Definition, Notation, Assumption
%% For proofs, please use the proof environment (the amsthm package is loaded by the MDPI class).

%=================================================================
% Full title of the paper (Capitalized)
\Title{Clearance-Based Labels for Aircraft Encounters: Why the Unit of Evaluation Matters and How to Check It}

% Author Orchid ID: enter ID or remove command
\newcommand{\orcidauthorA}{0000-0002-9097-1323} % Add \orcidA{} behind the author's name
\newcommand{\orcidauthorB}{0000-0002-0567-9572} % Add \orcidB{} behind the author's name
\newcommand{\orcidauthorC}{0000-0003-2258-6623} % Add \orcidC{} behind the author's name

% Authors, for the paper (add full first names)
\Author{Borja Rodríguez \textsuperscript{1,*}\orcidA{}, Félix Hernández-del-Olmo \textsuperscript{2}\orcidB{} and Elena Gaudioso \textsuperscript{2}\orcidC{}}


%\longauthorlist{yes}

% MDPI internal command: Authors, for metadata in PDF
\AuthorNames{Borja Rodríguez, Félix Hernández-del-Olmo and Elena Gaudioso}

% Affiliations / Addresses (Add [1] after \address if there is only one affiliation.)
\address{%
\textsuperscript{1} \quad Universidad Nacional de Educación a Distancia (UNED), Calle Bravo Murillo 38, 28015 Madrid, Spain; borja.rode@gmail.com\\
\textsuperscript{2} \quad Departamento de Inteligencia Artificial, Universidad Nacional de Educación a Distancia (UNED), Madrid, Spain; felixh@dia.uned.es (F.H.); elena@dia.uned.es (E.G.)}

% Corresponding author: personal, long-term email, as in the previous manuscript.
\corres{Correspondence: borja.rode@gmail.com}

% Current address and/or shared authorship
%\firstnote{Current address: Affiliation.}  
% Current address should not be the same as any items in the Affiliation section.

%\secondnote{These authors contributed equally to this work.}
% The commands \thirdnote{} till \eighthnote{} are available for further notes.

%\simplesumm{} % Simple summary

%\conference{} % An extended version of a conference paper

% Abstract (Do not insert blank lines, i.e. \\) 
\abstract{Historical air traffic control (ATC) clearances are increasingly used as labels of controller decisions, to learn conflict-resolution strategies and to calibrate encounter indicators. When such labels are built by marking every observation of an aircraft pair that is followed by a clearance within a short window, each encounter is counted as many times as it is observed, and its later observations, which are closer to the controller's action, outweigh its start; a clearance can also fall within several encounters. We show that this changes how strongly encounter geometry appears to predict controller action. Using the public Swedish Civil Air Traffic Control (SCAT) dataset, we build an encounter indicator from predicted minimum separation and the time to reach it, as a probe rather than a proposed detection method. Counted per observation, pairs predicted to lose separation within five minutes are three times as likely as the reference to be followed by a clearance; counted once per encounter, as a real-time tool would see them, about twice. The reduction concentrates in predicted losses of separation, appears with or without our calibrated levels, and is driven mainly by the repeated observation of the same encounter. Moreover, most of these encounters leave the recorded surveillance before their outcome can be checked. We propose reporting clearance-based labels with three checks: unit of evaluation, clearance multiplicity and surveillance coverage.}

% Keywords
\keyword{air traffic control; controller clearances; association strength; clearance attribution; encounter analysis; SCAT dataset} 

% The fields PACS, MSC, and JEL may be left empty or commented out if not applicable
%\PACS{J0101}
%\MSC{}
%\JEL{}

%%%%%%%%%%%%%%%%%%%%%%%%%%%%%%%%%%%%%%%%%%
% Only for the journal Diversity
%\LSID{\url{http://}}

%%%%%%%%%%%%%%%%%%%%%%%%%%%%%%%%%%%%%%%%%%
% Only for the journal Applied Sciences
%\featuredapplication{Authors are encouraged to provide a concise description of the specific application or a potential application of the work. This section is not mandatory.}
%%%%%%%%%%%%%%%%%%%%%%%%%%%%%%%%%%%%%%%%%%

%%%%%%%%%%%%%%%%%%%%%%%%%%%%%%%%%%%%%%%%%%
% Only for the journal Data
%\dataset{DOI number or link to the deposited data set if the data set is published separately. If the data set shall be published as a supplement to this paper, this field will be filled by the journal editors. In this case, please submit the data set as a supplement.}
%\datasetlicense{License under which the data set is made available (CC0, CC-BY, CC-BY-SA, CC-BY-NC, etc.)}

%%%%%%%%%%%%%%%%%%%%%%%%%%%%%%%%%%%%%%%%%%
% Only for the journal BioTech, Fishes, Neuroimaging and Toxins
%\keycontribution{The breakthroughs or highlights of the manuscript. Authors can write one or two sentences to describe the most important part of the paper.}

%%%%%%%%%%%%%%%%%%%%%%%%%%%%%%%%%%%%%%%%%%
% Only for the journal Encyclopedia
%\encyclopediadef{For entry manuscripts only: please provide a brief overview of the entry title instead of an abstract.}

%%%%%%%%%%%%%%%%%%%%%%%%%%%%%%%%%%%%%%%%%%
%% Different journals have different requirements. Please check the specific journal guidelines in the "Instructions for Authors" on the journal's official website.
%\addhighlights{%
%This is an optional section, the goal of which is to increase the discoverability and readability of the article via search engines and other scholars. Highlights should not be a copy of the abstract, but a simple overview allowing the reader to quickly and simply find out what the article is about and what can be cited from it. Each of these parts should be up to 2~bullet points.
%\minisec{What are the main findings?}
% \begin{itemize}[nosep]
% \item First bullet.
% \item Second bullet.
% \end{itemize}
%\minisec{What are the implications of the main findings?}
% \begin{itemize}[nosep]
% \item First bullet.
% \item Second bullet.
% \end{itemize}
%}

%%%%%%%%%%%%%%%%%%%%%%%%%%%%%%%%%%%%%%%%%%

%=================================================================
% Compatibility helpers for the figures and review markers in the original study.
% The MDPI class loads graphicx, caption, xcolor, amsmath and most table packages.
% MDPI uses \subfloat for multi-panel figures (see the official template).
% Do not load subcaption or the obsolete subfigure package simultaneously.
\usepackage{subfig}
\newlength{\figtw}
\AtBeginDocument{\setlength{\figtw}{\textwidth}}
\newcommand{\subimg}[2]{\parbox[c][#1\figtw][c]{\linewidth}{\centering\includegraphics[width=\linewidth,height=#1\figtw,keepaspectratio]{#2}}}
\newcommand{\tbc}[1]{\textcolor{red}{[TBC: #1]}}
\newcommand{\pend}{\textcolor{red}{pending}}
\newcommand{\figplaceholder}[1]{%
\fbox{\begin{minipage}[c][2.45cm][c]{0.92\linewidth}
\centering\sffamily\small\textbf{FIGURE PLACEHOLDER -- NOT GENERATED}\\[4pt]
\parbox{0.88\linewidth}{\centering #1}
\end{minipage}}}

\begin{document}

\section{Introduction}
\label{sec:intro}
Reliable identification of aircraft encounters that may require controller attention is central to air traffic control (ATC) and air traffic management (ATM). Tactical and medium-term conflict-detection tools rely on forecasts of aircraft position and separation, yet their operational value depends on more than geometric proximity. A forecast may not account for a planned level change, a previously issued clearance, traffic sequencing, or uncertainty in the flight trajectory. Conversely, the presence of a clearance cannot by itself establish that a specific pair of aircraft was in conflict. Understanding this distinction matters when historical controller actions are used as empirical targets for a detection system or as a reward signal for an intelligent decision-support agent.

Modern operational concepts already distinguish predicted losses of separation, encounters that constrain manoeuvring options, and groups of interactions that should be presented coherently to the controller. The medium-term conflict detection (MTCD) specification of EUROCONTROL, the European Organisation for the Safety of Air Navigation,\footnote{\url{https://www.eurocontrol.int}} explicitly recognises these functions, including uncertainty and the grouping of related encounters \cite{eurocontrol2017}. The research literature on conflict detection and resolution classifies methods by how they propagate predicted states and set detection thresholds \cite{kuchar2000}, and probabilistic formulations convert the growth of trajectory-prediction error with look-ahead time into a conflict probability \cite{paielli1997,prandini2000}. The separation minima that define a loss of separation are set by the procedures of the International Civil Aviation Organization (ICAO)\footnote{\url{https://www.icao.int}} \cite{icao4444}. In parallel, the availability of historical surveillance, flight plans and clearance logs has stimulated learning-based research on controller behaviour \cite{bastas2022,bastas2023,delgado2022}. These datasets provide a valuable observational foundation but rarely include an unambiguous label linking a particular clearance to the encounter that motivated it. In this paper, ``clearance'' is used in a broad sense for any controller input recorded in the system that instructs an aircraft, including what ICAO terminology calls instructions; only heading and level clearances are analysed.

The Swedish Civil Air Traffic Control (SCAT) dataset is especially attractive because it combines time-stamped surveillance with flight plans, system trajectory predictions and controller clearances for 13 sampled weeks of operations \cite{nilsson2023}. It has already been used to reconstruct potential conflicts resolved by heading or level clearances \cite{renkhoff2025}, to predict or reconstruct clearances with supervised learning \cite{renkhoff2026}, and to analyse pilot response times to clearances \cite{wustenbecker2025}. Thus, reconstructing encounters or detecting clearances from SCAT is not new in itself. The question explored here is how strongly pairwise geometry appears to predict clearance activity when every observation of an encounter is counted, compared with counting each encounter once, and how far the limits of later surveillance allow the outcome to be checked.

We restrict the study to upper en-route airspace: flight levels (FL; altitude in hundreds of feet on the standard pressure setting) FL290--FL410, that is, 29,000--41,000 ft or about 8.8--12.5 km, where aircraft are mostly in cruise. The restriction is deliberate. The labels examined here are meant to serve as the basis for subsequent work on learning-based conflict resolution for cruising traffic, for which the clearance patterns of climbing, descending and terminal-area traffic would be a different and confounding population.

\subsection{Related Work and Research Gap}\label{sec:related}
\subsubsection{Controller behaviour from operational data}
Bastas and Vouros \cite{bastas2022} formulated data-driven prediction of whether and when a controller reacts to a conflict, and subsequently studied models of the type of resolution action using historical flight information \cite{bastas2023}. Both contributions recognise the difficulty of recovering an operator's decision context from trajectories recorded after interventions have been applied. Delgado-Aguilera Jurado et al. \cite{delgado2022} evaluated potential conflict detection by controllers with a Bayesian network model calibrated from operational information. These studies establish the use of ATC events as proxies for detection or intervention. The present study builds on this direction and examines how the interpretation of such labels changes when every observation of the same encounter is counted as a separate example.

Other approaches seek controller-conformal conflict-resolution policies through imitation or behaviour cloning \cite{guleria2024}. Such models benefit from reliable examples of the situation that motivated an action. If each clearance is assigned to every pair containing the aircraft within a time window, the label may capture routine traffic management instead of resolution of the nominated pair. The central object of our study is therefore the integrity of the association between pairwise geometry and observed controller action, not the prediction of the specific action itself.

\subsubsection{Conflict reconstruction, threshold selection and prediction error}
Renkhoff \cite{renkhoff2025} reconstructed potential conflicts and compatible heading/level resolutions from the same SCAT source, while Renkhoff et al. \cite{renkhoff2026} demonstrated supervised learning for reconstructing clearance information. Gaume et al. \cite{gaume2025} extracted conflict-resolution situations from historical Automatic Dependent Surveillance--Broadcast (ADS-B) trajectories and estimated proximity-related thresholds using data. These are close precedents for the constituent operations of our pipeline. A useful distinction is that their scientific questions centre on recovering interventions or deconfliction situations, whereas ours examines how the association between predicted geometry and clearance activity depends on the unit of counting and on data coverage. Gaume et al. had earlier proposed a heuristic catalogue of resolved lateral conflicts from ADS-B and flight plans \cite{gaume2023}. In the SCAT-based study \cite{renkhoff2025}, the retained samples are those in which a clearance prevented a potential conflict (135 heading and 3,720 flight-level clearances); we instead keep all 51,950 eligible heading/level clearances and audit how many pair observations each clearance labels. Wüstenbecker et al. \cite{wustenbecker2025} analysed initiation delays and completion times of pilot responses on the same dataset (about 830,000 clearances), which bears directly on the choice of the 300-s association window examined in Section~\ref{sec:results_sensitivity} and Appendix~\ref{app:models}.

Xia et al. \cite{xia2024} investigated the interaction between trajectory predictor performance and the detection characteristics of adaptive conflict-detection systems. The dependence of alert quality on prediction uncertainty is therefore also established. Our descriptive comparison of predicted and observed separation adds a different kind of evidence because it is linked to identifiable clearance records and includes an explicit missing-surveillance category.

\subsubsection{Specific contribution being tested}
The literature above covers controller-action prediction, conflict reconstruction, threshold calibration, trajectory errors, and grouping of encounters. We have not identified, within these specific cited studies, an evaluation that examines \emph{jointly}, as its primary object, how the association between predicted geometry and clearance activity depends on the unit of counting (pair-instant versus episode), on the traceable many-to-many links between clearances and encounters, and on the observed trajectory coverage of the outcome. The calibrated four-level partition is only the vehicle for this audit, not the novelty. This gap is scoped to the studies cited above. Whether the reduction is specific to the four levels, and which part of the counting rule produces it, is examined in Section~\ref{sec:benchmarkresults}.

\subsection{Research Question and Contributions}
The practical question is how far the apparent strength of the association between encounter geometry and clearance-based labels depends on the unit at which the labels are counted. A label can be counted per \emph{pair-instant} (one aircraft pair at one evaluation time), per \emph{clearance} (one identified heading or level clearance) or per \emph{encounter episode} (a contiguous run of observations of the same two aircraft, Section~\ref{sec:encounters}). These units are related many-to-many: one aircraft pair appears at successive two-minute evaluations, and one clearance can label many pairs. Counting every observation of the same encounter as an independent example is analogous to temporal pseudoreplication \cite{hurlbert1984}: repeated measurements of one unit are treated as independent replicates. Treating these rows as independent decisions would overstate the amount of evidence, and an association that looks strong at one unit can look weaker at another. We therefore pose the following research question:
\begin{quote}
\emph{How much does counting every observation of an encounter, instead of the encounter, overstate how strongly predicted geometry is associated with controller clearances, why, and what should be reported so that clearance-based labels can be interpreted?}
\end{quote}

To answer it we use a case study on SCAT. Pair-instants with similar predicted geometry are grouped into four ordered levels, shown in green, yellow, orange and red (``non-green'' below means yellow, orange or red), according to the clearance activity that follows them. The grouping is a data-driven partition of geometry cells into groups that differ as much as possible in the subsequent clearance activity associated with them, and the number of groups is computed from the data rather than imposed (Section~\ref{sec:colors}). Green is the reference group, with no expected elevated activity, and each further level adds progressively higher associated activity. The levels use the minimum predicted normalised separation, $\mmin$, and the time to that minimum, $\tau$, defined in Section~\ref{sec:geometry}; their boundaries were selected from observed clearance associations rather than fixed by engineering judgement. The colours label ordered levels of expected ATC activity, not certified categories of incident severity.

The contribution of this paper is a quantified, reproducible audit of clearance-based labels, not a new conflict-detection or prediction method. Using SCAT as a case study, we measure how much the apparent strength of the association between predicted geometry and clearance activity changes with the unit of evaluation, why it changes, how many pair-instants and encounters a single clearance labels, and what share of encounters can be followed to an observable outcome. From this audit we derive three checks that we propose reporting with any clearance-based label (Table~\ref{tab:checks}): the unit of evaluation, clearance multiplicity and surveillance coverage. The audit reconstructs pair-instants, clearance events and encounter episodes, links them explicitly, and reports the denominators and exclusions of every step, so that each figure can be traced and reproduced from the public data and the released code. The four-level partition is used as a vehicle for the audit, not as a proposed standard, and we introduce no new geometric conflict metric, warning scheme or threshold-fitting procedure.

\begin{table}[t]
\caption{Proposed checks for clearance-based labels, how they are obtained in this study and the corresponding result on SCAT.}\label{tab:checks}
\small
\begin{tabular}{@{}p{0.15\linewidth}p{0.27\linewidth}p{0.22\linewidth}p{0.30\linewidth}@{}}
\toprule
Check & What to report & How it is obtained here & Result on SCAT\\
\midrule
Unit of evaluation & Label rate and lift by class per pair-instant and per episode, same window and green reference & Pair-instant and episode lifts under the same green definition and 300-s window (Tables~\ref{tab:units} and~\ref{tab:ladder}) & Red lift 3.11 per pair-instant versus 2.03 with the first observation of each encounter (1.94 with the episode's maximum colour; binary screen: 2.48 versus 1.80), mainly from repeated observation of the same encounter (Section~\ref{sec:benchmarkresults})\\
Clearance multiplicity & Pair-instants, distinct pairs, episodes and classes linked to each clearance & Many-to-many link table between clearances and pair-instants (Section~\ref{sec:attribution}) & Median of 2 pair-instants per linked clearance among non-green pairs (65 over all candidate pairs, an effect of the 350 NM radius); 27--41\% of clearances linked to non-green episodes relate to more than one episode (Section~\ref{sec:results_attribution})\\
Surveillance coverage & Share of encounters whose outcome can be observed over the full horizon, by class & Synchronised observed trajectories with at least 80\% grid coverage (Section~\ref{sec:observation}) & 28.8\% of red episodes are evaluable, versus 64.2\% of orange and 50.6\% of yellow (Section~\ref{sec:results_coverage})\\
\bottomrule
\end{tabular}
\end{table}

At the encounter level, the four-level partition (M0) is also compared with a geometric screen, continuous logistic models and monotone discrete alternatives (M1--M5) in Appendix~\ref{app:models}. A Gaussian conflict-probability model (M6) is deferred because its target differs from subsequent clearance activity.

The remainder of the paper is organised as follows. Section~\ref{sec:methods} describes the SCAT data, clearance eligibility, predicted geometry, calibration and retrospective auditing procedures, including the encounter-level model-comparison protocol. Section~\ref{sec:results} presents the empirical audit, the comparison of units and the sensitivity results, and Appendices~\ref{app:models} and~\ref{app:manoeuvres} give the model comparison and the manoeuvre screening. Section~\ref{sec:discussion} discusses their methodological and operational implications, and Section~\ref{sec:conclusion} concludes.

\section{Materials and Methods}\label{sec:methods}
\subsection{SCAT Dataset and Study Population}\label{sec:data}
SCAT comprises 13 non-consecutive weeks distributed from October 2016 to September 2017. The source describes approximately 170,000 scheduled flights, with military, private and incident-report-related flights removed \cite{nilsson2023}. The structural audit of our local archive counted 167,547 flights, 70,224,874 surveillance points, 1,187,437 raw clearance snapshots and 9,713,913 system-prediction updates. These are \emph{source-level} inventory counts, distinct from the much smaller set of flights and clearances eligible for the analysis.

The modelling study is restricted to the upper en-route airspace introduced in Section~\ref{sec:intro}, flight levels FL290--FL410 (29,000--41,000 ft), where aircraft are normally in cruise. At an evaluation instant, an unordered pair enters the candidate set when both aircraft are observed within this flight-level range and their current horizontal separation is below 350 nautical miles (NM; 1 NM = 1,852 m exactly, approximately 648 km). This wide initial radius deliberately preserves many geometrically unremarkable pairs as reference traffic. All eligible pairs are evaluated at 120-s intervals. Predictions extend 20 minutes into the future on a 10-s grid. The analysis therefore concerns normal operations as captured by SCAT, which excludes flights in incident reports.

\subsection{Clearance Provenance and Eligibility}
Clearance extraction records a new change in an actionable field rather than every repeated flight-plan snapshot. It yielded 1,019,799 distinct actionable records, of which 908,165 include a heading or level component. During preprocessing, we retained level clearances only when both assigned flight levels were numeric and different, both assigned levels and the recipient's observed flight level lay within FL290--FL410, and a valid surveillance state was available no more than 10 s before the clearance. Heading clearances required a valid numeric new heading that was either an initial assignment or a change from the previous assignment, together with a valid pre-clearance surveillance state within 10 s and an observed flight level in FL290--FL410. These rules selected 50,774 level and 1,176 heading records, giving 51,950 distinct \emph{eligible heading/level clearances}. The audit preserved each original clearance identifier and reconciled it with its flight, source member, archive week and coordinated universal time (UTC) epoch. Speed-only clearances were excluded from the target label.

The 51,950 selected records are a subset of the available heading/level changes. Results refer to this \emph{eligible subset}; a broader clearance definition is examined as a sensitivity analysis (Section~\ref{sec:results_sensitivity}).

\subsection{Temporal Partition and Observational Units}
Nine sampled weeks form the calibration partition and four form the evaluation partition. The four test-week start dates are 7 January, 1 April, 24 June and 16 September 2017. Colour selection was rerun using leave-one-week-out validation on the nine calibration weeks, and the resulting thresholds were held fixed during the final reported test computation. A previous exploratory version of the analysis had inspected all 13 weeks, so the four-week partition is \emph{held out from the final calibration run}, not from the whole study.

Three units are distinguished throughout the article. A \emph{pair-instant} is one unordered aircraft pair at one evaluation time. A \emph{clearance event} is one deduplicated, identified heading/level clearance. An \emph{encounter episode} is a contiguous run of observations for the same two aircraft, defined as described in Section~\ref{sec:encounters}. The relationships between the units are many-to-many and are retained throughout the audit.

\paragraph{Software.} The analysis was implemented in Python 3.14.6 using NumPy 2.3.5, pandas 3.0.3, SciPy 1.17.1, scikit-learn 1.9.0, Matplotlib 3.11.1 and PyArrow 25.0.1, among others.

\begin{figure}[tbp]
\centering
\includegraphics[width=\linewidth]{figures/fig01_workflow.png}
\caption{SCAT analysis workflow. Blue boxes denote observed source records and lavender boxes denote predictions, temporal links or derived summaries. The M0--M5 comparison is reported in Appendix~\ref{app:models}.}\label{fig:workflow}
\end{figure}

\subsection{Predicted encounter geometry}\label{sec:geometry}
For aircraft $i$ and $j$ evaluated at instant $t_0$, let $H_{ij}(t_0+s)$ and $V_{ij}(t_0+s)$ denote their predicted simultaneous horizontal and absolute vertical separations, respectively, at look-ahead $s$. The primary predictor follows the available horizontal route points at the current observed ground speed, follows the predicted altitude profile where available (with a constant-motion fallback when no usable route exists) and samples the 20-min forecast every 10 s. It is a simplified, route-informed predictor rather than a full four-dimensional model of planned turns, level changes or aircraft performance.

For the reference separation thresholds used in this analysis, we define
\begin{equation}\label{eq:geometry}
 m_{ij}(t_0+s)=\max\!\left\{\frac{H_{ij}(t_0+s)}{5\,\mathrm{NM}},\frac{V_{ij}(t_0+s)}{1000\,\mathrm{ft}}\right\},
\end{equation}
with predicted minimum and lead time
\begin{equation}\label{eq:minimum}
 \mmin(t_0)=\min_{s\in\{0,10,\ldots,1200\}\mathrm{s}}m_{ij}(t_0+s),\quad
 \tau(t_0)=\mathop{\mathrm{argmin}}_{s}m_{ij}(t_0+s).
\end{equation}
For this metric, $m<1$ means that the horizontal and vertical thresholds are both violated \emph{at the same predicted instant}; $m=2$ means that the larger normalised separation coordinate equals two, not that both distances simultaneously equal twice their thresholds. The limits (5 NM = 9,260 m and 1,000 ft = 304.8 m) are study parameters, not operational minima for every sector and route.

This maximum-coordinate definition produces nested boxes in the $(H,V)$ separation plane. It is transparent for the binary geometric screen $m<1$ and is used here as a descriptive indicator rather than as a calibrated probability. The predicted minimum and its occurrence time form a compact encounter descriptor whose association with clearance activity can be studied empirically.

\paragraph{Worked example.} Figure~\ref{fig:frame} shows these quantities at one evaluation instant (16 October 2016, 12:58 UTC) for a window of Swedish upper airspace containing eight aircraft between FL290 and FL410. Panel (a) shows, for each aircraft, its position, flight level, ground speed and vertical rate, together with its route-informed forecast for the next 20 minutes marked at 5, 10 and 20 minutes. The eight aircraft form 28 unordered pairs, and each pair is one pair-instant at this evaluation time. For every pair the forecast yields a time series of horizontal separation $H$ and vertical separation $V$ and hence of $m$; panel (b) shows it for the pair of aircraft A and B. Aircraft A climbs at about 1,700 ft/min (8.6 m/s) towards the level of aircraft B and the two fly almost head-on, so the vertical and horizontal thresholds are violated at the same time only around the crossing, where $m$ reaches $\mmin=0.39$ at $\tau=3.5$ min (210 s). Since $\mmin<1$ within five minutes, the pair meets the forecast-screening thresholds and is assigned to the highest intervention-activity level (red) defined in Section~\ref{sec:colors}. Of the 28 pairs, only this one has $\mmin<1$; one other pair comes close ($\mmin=1.66$) without violating the thresholds, and all remaining pairs keep $\mmin\ge 2$. This is the same encounter as in Fig.~\ref{fig:example}, where the forecast is compared with the later observed trajectories.

\begin{figure}[tbp]
\centering
\subfloat[Plan view (north up) of eight aircraft in upper en-route airspace, labelled by letter with flight level, ground speed and vertical rate.]{%
\begin{minipage}[t]{\linewidth}
\centering
\includegraphics[width=0.92\linewidth]{figures/fig10a_airspace_frame.png}
\label{framea}
\end{minipage}%
}\\[6pt]
\subfloat[Normalised separation of the pair A--B along the 20-min forecast.]{%
\begin{minipage}[t]{0.62\linewidth}
\centering
\includegraphics[width=\linewidth]{figures/fig10b_pair_geometry.png}
\label{frameb}
\end{minipage}%
}
\caption{Worked example of the encounter geometry at 12:58 UTC on 16 October 2016. (a) Dashed lines are route-informed 20-min forecasts, with markers at 5, 10 and 20 min; aircraft A and B (highlighted) are analysed in (b), and the black squares mark their positions at the closest approach. Base map: Natural Earth. (b) Horizontal ($H/5$ NM) and vertical ($V/1{,}000$ ft) components and their maximum $m$; the shaded band is $m<1$, and the circle marks $\mmin$ at lead time $\tau$.}\label{fig:frame}
\end{figure}

\subsection{Clearance association and traffic-density adjustment}\label{sec:lift}
At each pair-instant, the response $Y_{ij,t_0}=1$ if either member of the pair has at least one \emph{eligible} heading or level clearance in the interval $(t_0,t_0+300\,\mathrm{s}]$, and zero otherwise. This is a binary indicator of a temporally associated clearance; the event count remains at most one per pair-instant regardless of how many clearances occur. A clearance can label many pair-instants, including pairs that did not motivate it, so the label is an association and not a confirmed conflict.

Traffic is stratified using the concurrent number of aircraft inside FL290--FL410 over the study airspace as a whole. Calibration uses four strata, approximately 0--15, 15--25, 25--35 and more than 35 aircraft. This is the \emph{global instantaneous traffic count}, not a local sector density. Within each stratum $d$, let $r_{0d}$ be the observed clearance-label rate in a reference group (the pair-instants with no expected elevated activity, defined in Section~\ref{sec:colors}) estimated from calibration data. For a geometric group $g$ with $n_{gd}$ pair-instants, define its expected positives and stratified relative rate as
\begin{equation}\label{eq:lift}
 E_g=\sum_{d=1}^{4}n_{gd}r_{0d},\qquad L_g=\frac{O_g}{E_g},
\end{equation}
where $O_g$ is the number of pair-instants for which $Y=1$, \emph{not} the number of unique clearances. In the calibration split the reference group is fixed to $L=1$ by construction, because $r_{0d}$ is estimated from the same green pair-instants. The test-split lifts use density-specific green reference rates re-estimated on the four test weeks. Thus the lift of the reference group on the test weeks is 1 by construction; this is a descriptive held-out comparison rather than a prospective calculation with a frozen calibration reference rate. All lifts reported here use the green reference.

The observed stratified reference rates in calibration were approximately 6.9\%, 12.1\%, 13.3\%, and 13.3\% across the increasing traffic strata. Stratification adjusts for the overall traffic level only. The word ``expected'' in Eq.~\eqref{eq:lift} therefore denotes a \emph{standardised comparison} with the reference group, not a counterfactual estimate.

\subsection{Statistical calibration of discrete levels}\label{sec:colors}
Candidate geometric cells were formed from nine bands of $\mmin$ (below 1; 1--1.2; 1.2--1.5; 1.5--2; 2--3; 3--4; 4--6; 6--8; and 8--12) and four lead-time bands (below 2; 2--5; 5--10; and 10--20 min). Candidate intervention levels were fitted using the rate information within cells. The threshold-selection analysis compared constrained, operationally interpretable partitions with more flexible statistical groupings and considered different numbers $K$ of levels. For an ordered collection of cells, a profile-Poisson objective can be written, up to constants common to alternative partitions, as
\begin{equation}\label{eq:partition}
 \mathcal{Q}(\mathcal{P}_K)=\sum_{k=1}^{K}O_k\log\!\left(\frac{O_k}{E_k}\right),
\end{equation}
where $\mathcal{P}_K$ partitions the cells into $K$ contiguous groups in the chosen ordering; terms with $O_k=0$ are treated by continuity. A dynamic-programming segmentation step was used for that restricted search; exact segmentation by dynamic programming is reviewed in \cite{killick2012}. The operational colour scheme was selected by comparing predictive gains across nine leave-one-week-out calibration folds, with a simpler four-level specification preferred when the incremental improvement of a fifth level was small.

The frozen nominal rules used in the subsequent audits are
\begin{equation}\label{eq:color}
 C(\mmin,\tau)=
 \begin{cases}
 \mathrm{red},& \mmin<1,\ \tau<300\,\mathrm{s},\\
 \mathrm{orange},& \mmin<1,\ 300\le\tau\le1200\,\mathrm{s},\\
 \mathrm{yellow},&1.2\le\mmin<2,\\
 \mathrm{green},&\text{otherwise}.
 \end{cases}
\end{equation}
In particular, $1\le\mmin<1.2$ is green in this data-driven solution. The levels therefore describe the pattern of subsequent clearance activity under this forecasting and labelling pipeline, not four monotonically increasing \emph{geometric danger} classes.

Figure~\ref{fig:kselect} summarises how the number of levels was chosen. For each number of groups $K$, the score of a partition is its mean gain over a single group across the nine leave-one-week-out folds (the partition is fitted on eight weeks and scored on the ninth, with the five-minute label). The \emph{attainable} improvement is the gain of the best unconstrained partition of the candidate cells, which occurs at $K=11$ groups, and the \emph{share captured} by a scheme is its gain divided by this attainable gain. For unconstrained partitions the share is 68.9\%, 87.7\%, 92.2\%, 96.0\% and 97.3\% with two to six groups, and each further group adds about one percentage point or less (panel b). The operational four-level scheme of Eq.~\eqref{eq:color} captures 87.8\%; a three-level scheme (red for $\mmin<1$, yellow for $1.2\le\mmin<2$) captures 84.2\%; and the five-level scheme obtained by splitting the red level at two minutes captures 88.7\%, that is, 0.8 percentage points more than four levels. Four levels are therefore a compromise between captured score, interpretability and number of groups, not an optimum: the interpretable thresholds cost about four percentage points relative to the best unconstrained four-group partition, and a fifth level is a defensible alternative when a finer separation of the most active geometries is needed. Table~\ref{tab:increments} breaks the comparison down by step, using a paired $t$ statistic. Let $d_f$ be the difference between the held-out score of the richer and the simpler scheme in fold $f$ ($f=1,\dots,F$, with $F=9$ leave-one-week-out folds). The reported gain is the mean difference $\bar d$, its standard error is $\mathrm{SE}$, and
\begin{equation}\label{eq:tstat}
 t=\frac{\bar d}{\mathrm{SE}},\qquad \bar d=\frac1F\sum_{f=1}^{F}d_f,\qquad \mathrm{SE}=\frac{s_d}{\sqrt F},\qquad s_d^2=\frac1{F-1}\sum_{f=1}^{F}\left(d_f-\bar d\right)^2 ,
\end{equation}
which is Student's paired $t$ statistic. Under the null hypothesis of no average improvement and approximately normal fold differences, $t$ follows a Student distribution with $F-1=8$ degrees of freedom, whose two-sided 5\% critical value is 2.31; larger $|t|$ means that the gain is large relative to its fold-to-fold variability. Because the nine folds are consecutive weeks of one dataset and the schemes were chosen on the same data, $t$ is used as a descriptive measure of how consistent each step is, not as a formal significance test, and the number of folds improved (Folds in the table) is reported beside it. Separating red from orange by lead time is the only step beyond three levels that is clearly supported with the five-minute label ($t=5.2$, improvement in 8 of 9 folds); splitting the yellow level, either at $\mmin=1.5$ or by lead time, is not significant, so the fourth level comes from the red/orange split. A further split of red at two minutes is significant ($t=5.7$) but adds only about 1.6 units, around 0.8\% of the attainable 186.3. The conclusion depends on the clearance horizon: with two-minute clearances the red split at two minutes adds substantially ($t=7.7$), whereas with ten-minute clearances the red/orange split is marginal ($t=2.8$). Moving the lower yellow limit from 1.0 to 1.2 improves the score in all nine folds ($t=12.1$), which supports the green $1\le\mmin<1.2$ band of Eq.~\eqref{eq:color}. These percentages describe \emph{the segmentation criterion}, not the fraction of avoidable conflicts detected.

\begin{figure}[tbp]
\centering
\subfloat[Share of the attainable score captured by $K$ groups.]{%
\begin{minipage}[t]{0.49\linewidth}
\centering
\subimg{0.44}{figures/fig09a_k_share.png}
\label{kselecta}
\end{minipage}%
}\hfill
\subfloat[Gain from the $K$-th group (unconstrained partition).]{%
\begin{minipage}[t]{0.49\linewidth}
\centering
\subimg{0.44}{figures/fig09b_k_gain.png}
\label{kselectb}
\end{minipage}%
}
\caption{Choice of the number of levels (nine calibration weeks, leave-one-week-out, five-minute eligible-clearance label). In (a), the line is the unconstrained partition of the candidate cells into $K$ groups, the squares are the best operational scheme with $K$ levels, and the outlined square is the adopted four-level scheme. In (b), the gain is expressed in percentage points of the attainable score.}\label{fig:kselect}
\end{figure}

\begin{table}[tbp]
\caption{Incremental value of each step towards the interpretable schemes (nine calibration weeks, leave-one-week-out). Gain, $t$ and Folds refer to the five-minute label. Gain is the mean paired improvement of the profile-Poisson score (Eq.~\eqref{eq:partition}) on the held-out week with the five-minute label, $\pm$ one standard error across folds; $t$ is the paired statistic of Eq.~\eqref{eq:tstat}; ``Folds'' counts the weeks in which the step improved the score. The last two columns give $t$ when the label uses two- and ten-minute clearances. The final row compares two definitions of the yellow level.}\label{tab:increments}
\centering\small\setlength{\tabcolsep}{4pt}
\begin{tabular*}{\linewidth}{@{\extracolsep{\fill}}lrrrrr@{}}
\toprule
Step & Gain $\pm$ SE & $t$ & Folds & $t$ (2 min) & $t$ (10 min)\\
\hline
1 $\to$ 2 (add red/orange level) & $113.7\pm8.6$ & 13.2 & 9/9 & 19.9 & 10.3\\
2 $\to$ 3 (add yellow level) & $43.2\pm5.2$ & 8.2 & 9/9 & 8.0 & 6.6\\
3 $\to$ 4: split red/orange by lead time & $6.7\pm1.3$ & 5.2 & 8/9 & 1.8 & 2.8\\
3 $\to$ 4: split yellow at $\mmin=1.5$ & $0.1\pm0.4$ & 0.3 & 5/9 & $-0.5$ & 1.0\\
4 $\to$ 5: split red at 2 min & $1.6\pm0.3$ & 5.7 & 9/9 & 7.7 & $-0.7$\\
4 $\to$ 5: split yellow by lead time & $-0.2\pm0.2$ & $-1.4$ & 4/9 & 1.4 & 2.2\\
4 $\to$ 5: split yellow at $\mmin=1.5$ & $0.1\pm0.4$ & 0.3 & 5/9 & $-0.5$ & 1.0\\
Yellow from $\mmin\ge1.0$ to $\mmin\ge1.2$ & $25.4\pm2.1$ & 12.1 & 9/9 & 9.0 & 7.4\\
\bottomrule
\end{tabular*}
\end{table}


\begin{figure}[tbp]
\centering
\subfloat[Normalised-separation geometry with the region of each class shaded in its colour.]{%
\begin{minipage}[t]{0.42\linewidth}
\centering
\subimg{0.41}{figures/fig02a_geometry.png}
\label{colorsa}
\end{minipage}%
}\hfill
\subfloat[Clearance lift relative to the final green class.]{%
\begin{minipage}[t]{0.55\linewidth}
\centering
\subimg{0.41}{figures/fig02b_lift_heatmap.png}
\label{colorsb}
\end{minipage}%
}
\caption{Normalised-separation geometry and the four intervention-activity classes. In (b), calibration-week (nine weeks), traffic-density-adjusted lift of the five-minute eligible-clearance label relative to the final green class is shown; outlines mark the selected regions. The 1--1.2 band is green.}\label{fig:colors}
\end{figure}

\subsection{Traceable clearance-to-pair reconciliation}\label{sec:attribution}
The clearance-attribution audit replayed the pairwise prediction calculation for all 13 archive weeks and exported one row per eligible pair-instant to a canonical table. Each row preserves both aircraft identifiers, archive/source member, UTC evaluation time, $\mmin$, $\tau$, frozen colour and the original binary future-clearance label. A separate table retains the original clearance identifier and its timestamp. Attribution checks whether the clearance belongs to either aircraft in the pair and falls inside the stipulated future window. No deduplication is performed solely by matching coincident timestamps.

For each identified clearance $c$, let $A(c)$ be the set of labelled pair-instants containing its recipient. The multiplicity $|A(c)|$ is reported along with the number of distinct aircraft pairs and non-green episodes linked to that event. Reconciliation tests verify that the reconstructed binary flags reproduce the original binary future-clearance label for every pair-instant. This check verifies the internal consistency of the data pipeline, not the controller's motivation for a given clearance.

\subsection{Encounter aggregation and temporal clearance categories}\label{sec:encounters}
Pair-instants of the same two source aircraft are grouped into a single \emph{non-green episode} if consecutive non-green observations lie at most 360 s apart within the same archive week. The colour of an episode is the most severe colour it attained. Green pair-instants are not part of this non-green aggregation; the green episodes used in Section~\ref{sec:benchmarks} are built separately. A prior eligible clearance is a heading/level clearance to either member in the 600 s before episode onset; a later eligible clearance is one in the following 300 s of any pair-instant within the episode.

The episode categories \emph{prior only}, \emph{later only}, \emph{both}, and \emph{neither} are mutually exclusive after combining their two binary indicators. The window lengths are operational choices, not causal boundaries; their sensitivity is examined in Section~\ref{sec:results_sensitivity}. A flight can appear in several overlapping encounters and a clearance can be linked to several episodes, so grouping repeated observations reduces, but does not remove, the attribution ambiguity.

\subsection{Synchronized retrospective trajectory validation}\label{sec:observation}
For retrospective validation, we reconstruct the subsequent surveillance tracks of the two aircraft in an episode and evaluate their horizontal and vertical separation at the \emph{same} UTC grid instants. The forecast and observation are anchored to the first snapshot at which the episode attained its highest colour. The 20-min observed horizon uses a 10-s grid. Interpolation is allowed only across short gaps of at most 30 s, with no extrapolation beyond available surveillance points. An observed minimum is considered evaluable only when at least 80\% of the grid is covered and both endpoints are present. Asserting that an anticipated loss \emph{did not occur} additionally requires complete coverage.

An observed normalised minimum $m_{\mathrm{obs}}<1$ is a geometric screening outcome, not a certified separation infringement. Conversely, an observed minimum above one does not show that the original alert was unnecessary, because earlier clearances or route changes may have altered the trajectories. Episodes with insufficient follow-up are reported as \emph{not evaluable}, not relabelled as safe.

\begin{figure}[tbp]
\centering
\subfloat[Complete 20-minute horizon.]{%
\begin{minipage}[t]{0.49\linewidth}
\centering
\subimg{0.44}{figures/fig03a_full_horizon.png}
\label{examplea}
\end{minipage}%
}\hfill
\subfloat[Closest-approach detail.]{%
\begin{minipage}[t]{0.49\linewidth}
\centering
\subimg{0.44}{figures/fig03b_closest_approach.png}
\label{exampleb}
\end{minipage}%
}
\caption{Predicted and observed normalised separation for an example red encounter. The route-based forecast was frozen at the reference time (0 min); the two eligible clearances (numbered in (b)) and their estimated manoeuvre onsets are marked.}\label{fig:example}
\end{figure}

\subsection{Uncertainty and statistical dependence}\label{sec:uncertainty}
The calibration summary used 400 bootstrap replicates resampling whole days, with percentile-based 95\% intervals. This is preferable to treating millions of pair-instants as independent draws because repeated assessments of a flight and its neighbouring aircraft occur within a day. Day-level resampling does not remove systematic confounding or dependence across day boundaries. Because tests computed from overlapping cross-validation folds can be optimistic, the nine-fold paired $t$ statistics are used descriptively, and are complemented by fold-consistency summaries, week-level effect sizes and cluster-aware bootstrap resampling of whole days or of encounters \cite{field2007}.

\subsection{Encounter-Level Protocol}\label{sec:benchmarks}
The protocol is available in the code repository (\texttt{docs/decisiones\_protocolo.md}, release v1.1.0; see Data Availability). The test weeks were used repeatedly during implementation and sensitivity runs, so the results are exploratory held-out analyses rather than a single confirmatory test. The binary-screen comparison, the stepwise decomposition of the unit effect, the analysis by position within the episode and the robustness checks were added after the first results and are likewise exploratory.

\subsubsection{Estimands and cohorts}
\emph{Primary estimand.} For an encounter episode $e$ beginning at $T_0$, $Y^{\mathrm{ep}}_e=1$ if at least one eligible heading/level clearance to either aircraft occurs in $(T_0,T_0+300\,\mathrm{s}]$. This fixed window equalises exposure across episodes. The earlier outcome based on a 300-s window after \emph{any} episode observation is retained as an exposure sensitivity analysis. Prior-clearance status in the 600 s before $T_0$ is a covariate, not part of the outcome.

\emph{Cohorts.} The 54,441 non-green episodes of Section~\ref{sec:encounters} retain their identifiers and maximum colours. For the comparison of units (Tables~\ref{tab:units}--\ref{tab:position}), the reference is all green pair-instants, grouped into green-only runs under the same 360-s gap rule, so that every pair-instant of the test weeks belongs to exactly one episode (1,090,998 green and 17,338 non-green episodes; Table~\ref{tab:cohorts}). The model comparison of Appendix~\ref{app:models} instead uses only green observations with $\mmin<2$, excluding those within 360 s of a non-green observation for the same aircraft pair, which yields 12,926 green episodes and preserves 39,351 yellow, 9,039 orange and 6,051 red episodes. Results are reported under one-episode-one-vote and pair-instant weighting. A clearance linked to several episodes is kept in all of them; no ``true'' motivating pair is selected.

\begin{table}[t]
\caption{Cohorts used in the analyses on the four test weeks (the last row also gives the totals over the 13 weeks). The binary-screen groups depend on the threshold $c$; the values shown are for $c=1$, and the negative group includes the yellow class and the near-threshold green band.}\label{tab:cohorts}
\small\setlength{\tabcolsep}{4pt}
\begin{tabular}{@{}p{0.2\linewidth}p{0.26\linewidth}p{0.29\linewidth}p{0.15\linewidth}@{}}
\toprule
Analysis & Reference group & Compared group & Used in\\
\midrule
Pair-instants & 6,812,986 green pair-instants & 35,071 yellow, 6,218 orange and 2,605 red pair-instants & Tables~\ref{tab:pair} and~\ref{tab:units}\\
Episodes, four levels & 1,090,998 green episodes (all green pair-instants) & 17,338 non-green episodes: 12,548 yellow, 2,927 orange, 1,863 red & Tables~\ref{tab:units}, \ref{tab:ladder} and~\ref{tab:position}\\
Episodes, binary screen & 1,103,546 negative episodes & 4,790 positive episodes ($\mmin<1$) & Tables~\ref{tab:binary} and~\ref{tab:ladder}\\
Model comparison & 3,943 green episodes with $\mmin<2$ (12,926 in 13 weeks) & 17,338 non-green episodes (54,441 in 13 weeks) & Appendix~\ref{app:models}\\
\bottomrule
\end{tabular}
\end{table}

\emph{Splits.} The four test weeks are used for all comparisons of units; the thresholds of the four levels were calibrated on the nine calibration weeks (Section~\ref{sec:colors}).

\emph{Temporal interpretation.} Episode maximum colour and minimum geometry are defined using the full episode. They may contain observations after the start-based window or after a clearance has occurred. Thus this is a held-out \emph{retrospective association} comparison, not a deployable $T_0$ predictor.

\subsubsection{Sensitivity Analyses}
The benchmark sensitivity analysis changes the clearance-association horizon (including a fixed 120-s and 600-s window and the original any-observation exposure rule), the prior-clearance horizon, the episode-gap threshold, the inclusion limit for green episodes, and the clearance-eligibility definition. These variants retain the same distinction between calibration and evaluation weeks. The complete quantitative findings are reported in Section~\ref{sec:results_sensitivity}.

\subsubsection{Clearance Eligibility and Surveillance-Coverage Audit}
The audit traces the successive eligibility filters for extracted heading and level clearances, avoiding interpretation of overlapping exclusion criteria as disjoint causes. For episodes without evaluable follow-up, the analysis examines individual-aircraft native tracks at the 20-min endpoint and distinguishes out-of-band last observations, in-band track endings and unavailable start states. These classifications describe available evidence rather than prove operational causes of track loss. Quantitative findings are reported in Section~\ref{sec:results_sensitivity}.

\section{Results}\label{sec:results}
\subsection{Pair-instant classification and week-level evaluation}
The canonical reconstruction contains 21,567,488 pair-instants over the 91 nominal SCAT days. The event-level reconciliation reproduces the original binary future-clearance labels, with zero disagreements. Table~\ref{tab:pair} reports exact counts and final green-referenced lifts. These are rates of \emph{pair-instants with a later eligible clearance}, not rates of independent conflicts or distinct controller decisions. The calibrated rate increases from green through red.

\begin{table}[t]
\caption{Frozen four-class summary for the nine calibration weeks and four evaluation weeks. ``With label'' denotes an eligible heading/level clearance within 300 s of that pair-instant; rows are not independent events. Confidence intervals are the reported day-bootstrap 95\% intervals. Test-split lifts use green reference rates re-estimated on the test weeks.}\label{tab:pair}
\begin{tabular}{@{}llrrrl@{}}
\toprule
Split & Class & Pair-instants & With label & Rate (\%) & Lift (95\% confidence interval)\\
\midrule
Cal. & Green & 14,616,594 & 1,906,874 & 13.0 & 1.00 (ref.)\\
Cal. & Yellow & 75,289 & 15,361 & 20.4 & 1.56\\
Cal. & Orange & 12,896 & 3,641 & 28.2 & 2.15\\
Cal. & Red & 5,829 & 2,223 & 38.1 & 2.91\\
Test & Green & 6,812,986 & 919,826 & 13.5 & 1.00 (ref.)\\
Test & Yellow & 35,071 & 7,269 & 20.7 & 1.53 (1.48--1.59)\\
Test & Orange & 6,218 & 1,874 & 30.1 & 2.23 (2.10--2.34)\\
Test & Red & 2,605 & 1,100 & 42.2 & 3.11 (2.91--3.32)\\
\bottomrule
\end{tabular}
\end{table}

The relative differences are consistent with a graded intervention-activity response under the chosen denominator. The 2,605 red pair-instants in test are repeated observations, not 2,605 independent encounters: the same aircraft pair can be sampled repeatedly, and a clearance unrelated to the pair's geometry can produce a positive label.

\begin{figure}[tbp]
\centering
\subfloat[Label rate (pair-instant denominators above bars).]{%
\begin{minipage}[t]{0.49\linewidth}
\centering
\subimg{0.43}{figures/fig04a_label_rate.png}
\label{lifta}
\end{minipage}%
}\hfill
\subfloat[Relative lift versus green.]{%
\begin{minipage}[t]{0.49\linewidth}
\centering
\subimg{0.43}{figures/fig04b_relative_lift.png}
\label{liftb}
\end{minipage}%
}
\caption{Eligible clearance-label rates for pair-instants in calibration and test, with 95\% intervals from 400 whole-day bootstrap replicates (a), and density-adjusted green-referenced lifts (b). Numbers above bars give pair-instant denominators. Bars count repeated pair-instants with a later clearance, not independent encounters.}\label{fig:lift}
\end{figure}

\subsection{Multiplicity of clearance attribution}\label{sec:results_attribution}
Of 51,950 eligible unique clearances, 47,339 were associated with at least one pair-instant by the specified five-minute matching rule; the remaining 4,611 had no linked pair-instant and are not shown in Fig.~\ref{fig:multiplicity}(a). In total, 3,176,638 clearance-to-pair-instant links were found, exceeding the 2,858,168 pair-instants with a positive label because some pair-instants had more than one clearance in their window. The median number of linked pair-instants per linked clearance was 65, with 90th percentile 108. This count depends on the 350 NM candidate radius, which includes many geometrically unremarkable neighbours: the median falls to 22 pair-instants within 100 NM and to 7 within 50 NM. Restricted to non-green pair-instants it is 2 (90th percentile 4; 18,668 clearances with at least one link) and restricted to $\mmin<1$ it is 1 (6,939 clearances); the median number of distinct aircraft pairs is 33 for all candidate pairs and 1 for non-green pairs. Among relevant encounters, the multiplicity per clearance in pair-instants is therefore small, and what persists is the link of one clearance to several episodes, described below.

Importantly, a distinct clearance may be counted in more than one colour if the receiving aircraft participates in different nearby pairs: 6,383 of the 28,324 distinct linked clearances appear in more than one colour. Similarly, at episode level there were 45,463 episode--clearance links representing only 28,324 distinct linked clearances, of which 11,474 were linked to more than one episode. Counts of ``clearances in red'' therefore depend on the unit: 45,463 episode--clearance links, 35,102 clearances unique within each colour, or 28,324 clearances unique across all colours. Any paper reporting such counts should state which one is used. The share of linked clearances that belong to more than one episode decreases with the radius used to build the episodes but does not disappear: 40.5\% at 350 NM, 32.8\% (8,015 of 24,405) at 100 NM and 27.0\% (5,587 of 20,680) at 50 NM.

\begin{figure}[tbp]
\centering
\subfloat[Linked pair-instants per linked clearance.]{%
\begin{minipage}[t]{0.33\linewidth}
\centering
\subimg{0.33}{figures/fig05a_links_per_order.png}
\label{multiplicitya}
\end{minipage}%
}\hfill
\subfloat[Distinct flight pairs per linked clearance.]{%
\begin{minipage}[t]{0.33\linewidth}
\centering
\subimg{0.33}{figures/fig05b_pairs_per_order.png}
\label{multiplicityb}
\end{minipage}%
}\hfill
\subfloat[Event and link totals.]{%
\begin{minipage}[t]{0.28\linewidth}
\centering
\subimg{0.33}{figures/fig05c_link_totals.png}
\label{multiplicityc}
\end{minipage}%
}
\caption{Many-to-many association of eligible controller clearances with pairwise observations. Panel (a) shows only the 47,339 clearances with at least one linked pair-instant; the remaining 4,611 eligible clearances have none. One clearance may label several pairs and repeated snapshots. This is label multiplicity, not controller workload.}\label{fig:multiplicity}
\end{figure}

\subsection{Non-green encounter episodes and clearance timing}
The encounter aggregation identifies 54,441 non-green episodes across the 13 weeks: 39,351 yellow, 9,039 orange and 6,051 red under their maximum attained colour. Within the four evaluation weeks there were 1,863 red episodes. Their mutually exclusive clearance-timing categories were 642 later-only, 213 both prior and later, 187 prior-only, and 821 with neither eligible clearance in the selected windows. Hence, 855/1,863 (45.9\%) had a later eligible clearance and 400/1,863 (21.5\%) had a prior eligible clearance, while 821/1,863 (44.1\%) had neither; these overlapping marginals are not additive. The 42.2\% rate reported for red \emph{pair-instants} is therefore not the same estimand as the 45.9\% share of red \emph{episodes} with a later eligible clearance.

These values depend on the episode construction, eligibility filters and windows. They do not indicate whether the 821 red encounters without a clearance in the windows were resolved earlier, followed a planned manoeuvre or needed an unrecorded response. A subsequent broader examination found that 99 of these 821 episodes had an eligible clearance 10--30 min before onset, and 574 had \emph{some} recorded clearance in the preceding 30 min. These broader categories are reported descriptively and are not interpreted as confirmed resolutions.

\begin{figure}[tbp]
\centering
\includegraphics[width=\linewidth]{figures/fig06_red_episode_orders.png}
\caption{Mutually exclusive timing categories for 1,863 red encounter episodes in the four test weeks. Prior means an eligible clearance in the 600 s before episode onset; later means an eligible clearance within 300 s after any episode observation. One clearance may link to several encounters. The 42.2\% rate (1,100/2,605 red pair-instants) and the 45.9\% share (855/1,863 red episodes with a later clearance) use different denominators.}\label{fig:episodeorders}
\end{figure}

\subsection{Follow-up coverage and observed separation}\label{sec:results_coverage}
A major result is the limitation imposed by subsequent surveillance. Of 54,441 non-green episodes, 27,451 (50.4\%) met the specified criterion for an evaluable synchronised observed minimum; 26,990 were excluded for insufficient follow-up. Of the evaluable set, 27,447 had complete 10-s grid coverage. Table~\ref{tab:coverage} shows that the evaluability rate depends strongly on colour: only 1,744/6,051 red episodes (28.8\%) are evaluable, versus 5,802/9,039 orange (64.2\%) and 19,905/39,351 yellow (50.6\%). Since the opportunity to observe the subsequent minimum is not uniform, evaluating only available episodes can distort comparisons between colours.

\begin{table}[t]
\caption{Encounter-level coverage and geometric screening outcome across all 13 SCAT weeks. ``Observed below 1'' means a synchronised measured value of $m<1$, not a certified operational incident.}\label{tab:coverage}
\begin{tabular}{@{}lrrrrr@{}}
\toprule
Class & Episodes & Evaluable & Evaluable (\%) & Obs. $m<1$ & Pred. loss not obs.\\
\midrule
Yellow & 39,351 & 19,905 & 50.6 & 349 & Not applicable\\
Orange & 9,039 & 5,802 & 64.2 & 470 & 5,332\\
Red & 6,051 & 1,744 & 28.8 & 703 & 1,040\\
\midrule
Total & 54,441 & 27,451 & 50.4 & 1,522 & --\\
\bottomrule
\end{tabular}
\end{table}

For red encounters, 703 evaluable episodes had an observed normalised minimum below one; 1,040 had a forecast loss not observed under complete follow-up (one additional evaluable episode did not meet that stricter complete-coverage criterion). These counts describe the disagreement between a simplified forecast and the later observed trajectories; because manoeuvres may have changed the trajectories, they are not detector precision or false-alarm rates.

In the four evaluation weeks specifically, 821 red episodes had no eligible clearance in the original timing windows. Only 150 of those were evaluable; 89 had observed $m<1$ and 61 had complete follow-up with observed $m\ge1$. Of those 61, 46 had at least one recorded clearance in a broadened 30-min preceding window and 13 had none in that broadened window or the episode/five-minute follow-up; the remaining 2 cases had no recorded clearance in the preceding 30 minutes but had at least one during the episode/five-minute follow-up. Their median normalised prediction discrepancy was approximately 0.020, and only one exceeded 0.1. These residual cases are not explained by the clearance-linking procedure and are not interpreted as false alarms.

\begin{figure}[tbp]
\centering
\subfloat[Surveillance coverage.]{%
\begin{minipage}[t]{0.49\linewidth}
\centering
\subimg{0.39}{figures/fig07a_coverage.png}
\label{coveragea}
\end{minipage}%
}\hfill
\subfloat[Observed geometric screen.]{%
\begin{minipage}[t]{0.49\linewidth}
\centering
\subimg{0.39}{figures/fig07b_observed_screen.png}
\label{coverageb}
\end{minipage}%
}
\caption{Surveillance evaluability and observed geometric screening across all 13 weeks. Evaluable requires at least 80\% of the 10-s grid and both endpoints; absence of an observed predicted loss is claimed only under complete coverage, and one red evaluable episode has partial coverage. Observed $m<1$ is a geometric screen, not a certified incident.}\label{fig:coverage}
\end{figure}

\subsection{Comparison of Units: Pair-Instants, Episodes and Repeated Observation}\label{sec:benchmarkresults}
Table~\ref{tab:units} compares the units on the four test weeks, with the green class unrestricted in $\mmin$ as at the pair-instant level (cohort of Table~\ref{tab:cohorts}), the same 300-s later window and lifts standardised by the instantaneous aircraft count. The gradient is present at both units but weaker per episode: the red lift is 1.94 per episode against 3.11 per pair-instant, with non-overlapping confidence intervals. This is the lift with the maximum colour of the episode; keeping only its first observation, which is what a real-time tool would see, gives 2.03 (Table~\ref{tab:ladder}). Anchoring the window as at the pair-instant level, that is, 300 s after any observation of the episode, lowers the lifts to 0.88, 1.18 and 1.32, because the green rate rises to 34.7\%, so yellow falls below green. Part of this effect is mechanical, because a longer episode has more windows and is more likely to contain a clearance; it illustrates why the exposure window has to be fixed when episodes are compared.

\begin{table}[t]
\caption{Test-week lifts of the eligible-clearance rate relative to green, by unit of evaluation and window anchor (cohort of Table~\ref{tab:cohorts}). Green is unrestricted in $\mmin$ in all columns, the later window is 300 s and lifts are standardised by the instantaneous aircraft count; 95\% day-bootstrap intervals in parentheses. Rates in \%.}\label{tab:units}
\small
\begin{tabular}{@{}lrrrrrrrr@{}}
\toprule
 & \multicolumn{2}{c}{Pair-instant} & & \multicolumn{2}{c}{Episode, from start} & & \multicolumn{2}{c}{Episode, any observation}\\
\cmidrule(lr){2-3}\cmidrule(lr){5-6}\cmidrule(lr){8-9}
Class & Rate & Lift & & Rate & Lift & & Rate & Lift\\
\midrule
Green & 13.5 & 1.00 (ref.) & & 15.8 & 1.00 (ref.) & & 34.7 & 1.00 (ref.)\\
Yellow & 20.7 & 1.53 (1.49--1.59) & & 21.5 & 1.36 (1.30--1.42) & & 30.5 & 0.88 (0.85--0.91)\\
Orange & 30.1 & 2.23 (2.11--2.34) & & 27.7 & 1.74 (1.64--1.83) & & 41.0 & 1.18 (1.14--1.23)\\
Red & 42.2 & 3.11 (2.92--3.32) & & 30.9 & 1.94 (1.83--2.06) & & 45.9 & 1.32 (1.27--1.37)\\
\bottomrule
\end{tabular}
\end{table}

To check that this pattern does not depend on the four calibrated levels, the comparison was repeated with the binary screen M1, positive when $\mmin<c$ and negative otherwise (the negative group includes yellow and the near-threshold green band), for $c=1$ (the screen M1 of the model comparison), 1.5 and 2 (Table~\ref{tab:binary}). Since $\mmin<1$ is exactly the union of the red and orange levels, this shows that the effect does not depend on the calibration or on the cut by lead time, not that it does not depend on the geometry. For $c=1$ the positive group has 8,823 pair-instants, of which 2,974 (33.7\%) are followed by an eligible clearance, against 13.5\% among the 6,848,057 negatives; the standardised lift is 2.48 per pair-instant and 1.81 per episode when the window starts with the episode, a ratio of 0.73 (95\% interval 0.70--0.76), and 1.24 when the window is anchored at any observation (ratio 0.50). The reduction from pair-instants to episodes is therefore not specific to the four levels, and it is concentrated in predicted losses of separation: for the looser thresholds the pair-instant lift is already only 1.56 and 1.54, and the ratio with the window from the episode start is close to one (0.97 for $c=1.5$ and 0.94 for $c=2$), whereas anchoring the window at any observation lowers the lift for all three thresholds (ratios 0.50, 0.81 and 0.72). The ratios are computed within each bootstrap replicate, with the same resampled days for both units.

\begin{table}[t]
\caption{Standardised lift of the binary screen $\mmin<c$ relative to the negative group, test weeks, 300-s later window (cohort of Table~\ref{tab:cohorts}). Ratio is the episode lift divided by the pair-instant lift, with 95\% day-bootstrap intervals computed within each replicate.}\label{tab:binary}
\centering\small\setlength{\tabcolsep}{4pt}
\begin{tabular}{@{}crrlrl@{}}
\toprule
 & Pair-instant & \multicolumn{2}{c}{Episode, from start} & \multicolumn{2}{c}{Episode, any observation}\\
\cmidrule(lr){3-4}\cmidrule(lr){5-6}
$c$ & Lift & Lift & Ratio & Lift & Ratio\\
\midrule
1 & 2.48 & 1.81 & 0.73 (0.70--0.76) & 1.24 & 0.50 (0.48--0.51)\\
1.5 & 1.56 & 1.51 & 0.97 (0.94--1.00) & 1.27 & 0.81 (0.80--0.83)\\
2 & 1.54 & 1.45 & 0.94 (0.92--0.96) & 1.11 & 0.72 (0.71--0.73)\\
\bottomrule
\end{tabular}
\end{table}

To separate what changes between the pair-instant and the episode, the comparison was decomposed into a sequence of analytic definitions (Table~\ref{tab:ladder}, Fig.~\ref{fig:ladder}), always with the same 300-s window length, the green (or negative) reference and the traffic stratification by instantaneous aircraft count. Step 0 is the pair-instant analysis. Step 0b keeps exactly the same pair-instants, colours, outcomes, reference group and traffic strata as Step 0 and only weights each observation by $1/n_e$, where $n_e$ is the number of observations of its episode, so that each episode contributes the same total weight (the total weight equals the number of episodes). Step 1 keeps only the first observation of each episode. Step 2 assigns to the episode its maximum colour and opens the window at the start of the episode, which is the definition of Table~\ref{tab:units}. Step 3 keeps the same episode colour but opens the window at the first instant at which the maximum colour is reached ($T_{\max}$); since $T_{\max}$ is known only once the episode has been observed, this step is retrospective.

The red lift falls from 3.11 (Step 0) to 2.66 when only the weighting changes (ratio 0.86, 95\% interval 0.83--0.88), to 2.03 when only the first observation is kept (0.65; 0.62--0.69) and to 1.94 with the episode colour (0.63; 0.59--0.66). Between consecutive steps, the paired ratios are 0.76 (0.74--0.79) from 0b to 1, a large change, but only 0.96 (0.93--0.99) from 1 to 2, so assigning the maximum colour of the episode instead of the initial one changes the red lift little; for the binary screen the corresponding ratios are 0.80 (0.78--0.82) and 1.01 (0.98--1.03). Anchoring the window at $T_{\max}$ raises it to 2.28, a factor of 1.18 (1.12--1.24) above Step 2, but it remains 0.73 (0.70--0.77) of the pair-instant value. The maximum colour is reached after the first observation in only 1,206 of the 17,338 non-green test episodes (7\%), so the anchoring changes the window for a small share of episodes; the red class gains because those episodes are, by construction, observed in a state of greater predicted proximity. For the binary screen $\mmin<1$ the decrease is similar (2.48, 2.26, 1.80 and 1.81) and anchoring at $T_{\max}$ changes the lift very little (1.85; ratio 1.02 relative to Step 2; overall ratio 0.75). The reduction therefore persists once the recount of observations, the assignment of the episode colour and the anchoring of the window have been controlled for in turn, although the red lift in this decomposition is smaller than for the pair-instant under every definition. Replacing the stratifier by the number of episode starts per hour changes the red lifts by at most 0.01. These steps are differences between analytic definitions, not an exact causal decomposition, and they were added after the first results (Section~\ref{sec:benchmarks}). All steps use the cohort of Table~\ref{tab:cohorts} (1,108,336 test episodes), in which every pair-instant belongs to exactly one episode.

\begin{table}[t]
\caption{Decomposition of the pair-instant-to-episode reduction: standardised lift versus green (four levels) or versus the negative group (binary screen $\mmin<1$), test weeks, 300-s later window, stratified by instantaneous aircraft count. 95\% day-bootstrap intervals in parentheses. The last row is retrospective.}\label{tab:ladder}
\centering\footnotesize\setlength{\tabcolsep}{2.5pt}
\begin{tabular}{@{}llccc|c@{}}
\toprule
Step & Unit and window & Yellow & Orange & Red & $\mmin<1$\\
\midrule
0 & Pair-instant & 1.53 (1.49--1.59) & 2.23 (2.11--2.34) & 3.11 (2.92--3.32) & 2.48 (2.36--2.60)\\
0b & Pair-instant, weight $1/n_e$ & 1.69 (1.63--1.76) & 2.06 (1.94--2.18) & 2.66 (2.52--2.82) & 2.26 (2.14--2.36)\\
1 & First observation & 1.39 (1.34--1.45) & 1.71 (1.62--1.80) & 2.03 (1.91--2.16) & 1.80 (1.72--1.89)\\
2 & Episode, window from start & 1.36 (1.30--1.42) & 1.74 (1.64--1.83) & 1.94 (1.83--2.06) & 1.81 (1.74--1.89)\\
2$'$ & Episode, any observation & 0.88 (0.85--0.91) & 1.18 (1.14--1.23) & 1.32 (1.27--1.37) & --\\
3 & Episode, window from $T_{\max}$ & 1.36 (1.30--1.42) & 1.74 (1.64--1.83) & 2.28 (2.16--2.42) & 1.85 (1.77--1.93)\\
\bottomrule
\end{tabular}
\end{table}

\begin{figure}[tbp]
\centering
\subfloat[Four levels, lift versus green.]{%
\begin{minipage}[t]{0.49\linewidth}
\centering
\subimg{0.41}{figures/fig11a_ladder_four_class.png}
\end{minipage}%
}\hfill
\subfloat[Binary screen $\mmin<1$, lift versus negative.]{%
\begin{minipage}[t]{0.49\linewidth}
\centering
\subimg{0.41}{figures/fig11b_ladder_binary.png}
\end{minipage}%
}
\caption{Lift along the sequence of definitions of Table~\ref{tab:ladder} (Steps 0, 0b, 1, 2 and 3), with 95\% day-bootstrap intervals. Step 3 anchors the window at the first instant of the maximum colour and is retrospective.}\label{fig:ladder}
\end{figure}

The decomposition places the largest step between Step 0b and Step 1 (red ratio 0.76), that is, in which observations of an encounter are counted. Table~\ref{tab:position} shows what differs between them: the label rate rises with the position of the observation within its episode, from 32.3\% for the first observation of a red episode to 46.9\%, 53.1\% and 57.1\% for the second, the third and the fourth or later, and from 28.8\% to 42.8\% for $\mmin<1$. This is consistent with an encounter that persists until the controller acts, or with the time $\tau$ to the predicted minimum shortening as the encounter unfolds, since controllers may act when the crossing is closer: later observations are closer to the action, and counting each of them as an independent example raises the pair-instant rate above the rate of the encounter at its start (red: 42.2\% against 30.9\%). The reference also changes with the unit: the green rate is 13.5\% per observation and 15.8\% per episode, which lowers the lift further. Because the label looks 300 s ahead at 120-s steps, one clearance can mark at most three consecutive observations of the same pair; what differs between units is which observations enter the denominator, not how many times a clearance is counted.

\begin{table}[t]
\caption{Label rate and lift by position of the observation within its episode, test weeks, for the red level and for the binary screen $\mmin<1$. Rate is the share of observations followed by an eligible clearance within 300 s; the lift uses the green (negative) pair-instant reference of Step 0 of Table~\ref{tab:ladder}, stratified by the instantaneous aircraft count, with 95\% day-bootstrap intervals.}\label{tab:position}
\centering\footnotesize\setlength{\tabcolsep}{3pt}
\begin{tabular}{@{}lrrcrrc@{}}
\toprule
 & \multicolumn{3}{c}{Red} & \multicolumn{3}{c}{$\mmin<1$}\\
\cmidrule(lr){2-4}\cmidrule(lr){5-7}
Position & Obs. & Rate (\%) & Lift & Obs. & Rate (\%) & Lift\\
\midrule
1 & 1,274 & 32.3 & 2.38 (2.22--2.54) & 4,003 & 28.8 & 2.11 (2.01--2.23)\\
2 & 590 & 46.9 & 3.46 (3.10--3.88) & 2,077 & 34.2 & 2.52 (2.40--2.65)\\
3 & 275 & 53.1 & 3.89 (3.48--4.27) & 1,117 & 37.2 & 2.74 (2.48--3.01)\\
$\geq 4$ & 466 & 57.1 & 4.19 (3.74--4.69) & 1,626 & 42.8 & 3.15 (2.91--3.44)\\
\bottomrule
\end{tabular}
\end{table}

The ratios between units are stable. With episode gaps of 240 and 480 s the paired ratio of the first-observation lift to the pair-instant lift for the red level is 0.65 (0.61--0.69) and 0.65 (0.62--0.69), as with 360 s (0.65), and the ratio of the episode lift to the pair-instant lift for $\mmin<1$ is 0.73 (0.70--0.77) and 0.73 (0.70--0.76). In the nine calibration weeks, on which the colour thresholds were fitted, the corresponding ratios are 0.68 (0.64--0.71) for red and 0.73 (0.71--0.75) for $\mmin<1$.

\subsection{Sensitivity and Coverage Audits}\label{sec:results_sensitivity}
\subsubsection{Association Windows}
The first exposure sensitivity compares the fixed start-based clearance window with a 300-s window after any observation in the episode (Table~\ref{tab:units}). At a fixed 120-s later window, the red clearance rate is 13.7\%, below orange at 14.4\%; the four-class monotone ordering therefore does not persist at this shorter horizon. At 600 s, class-rate ordering is restored. With 240- and 480-s episode gaps, the class-rate ordering is retained. Relabelling the same episodes with all 908,165 extracted heading/level records raises clearance rates sharply: in test they are 23.6\%, 36.7\%, 58.3\% and 34.5\% for green through red, so red now falls below orange. The sensitivity of the model comparison to these choices is reported in Appendix~\ref{app:models}.

\subsubsection{Clearance Eligibility and Surveillance Missingness}
The source extraction contains 1,019,799 actionable rows, including 908,165 heading/level records. The frozen eligible set comprises 50,774 level and 1,176 heading clearances (51,950 total). For level records, 731,166 candidates fall to 563,511 with two numeric authorised levels, 416,449 with a changed flight level, 360,475 with a valid previous focal state, 113,163 in the analysed flight-level band and 50,774 finally eligible. Heading clearances follow a separately audited cascade; its 177,004 candidate rows yield 1,176 eligible events in the band. Among the broader heading/level records excluded from the eligible set, 543,535 have a matched flight level below FL290 and 248,889 lack one; 63,757 are recorded in FL290--FL410 but fail another eligibility condition. The steps are overlapping criteria, not mutually exclusive causes.

The observed-track screen makes 19,905/39,351 yellow, 5,802/9,039 orange and 1,744/6,051 red episodes evaluable. A second pass over native individual tracks partitions all 26,990 unevaluable episodes by observable evidence at the 20-min endpoint. In the red class, 479 have at least one aircraft last observed outside FL290--FL410, 3,823 have a track ending before the endpoint while last observed inside that band, and five lack an observed start state. In the yellow and orange classes, the corresponding in-band track-end counts are 14,479 and 1,089. Three yellow episodes reach the nominal archive boundary. No episode in this hierarchy has a later track that brackets the missing endpoint across a gap. An in-band track end is consistent with leaving the recorded Swedish surveillance domain: the last position of 98.9\% of the red in-band track ends lies within 12 NM of the empirical edge of the recorded data (Section~\ref{sec:limits}), although no operational airspace polygon is defined. These are evidence categories, not causal diagnoses of missingness; no absence-of-conflict conclusion is drawn for unevaluable episodes.

\section{Discussion}\label{sec:discussion}
\subsection{Why the observational unit matters}
The observed gradient of clearance-label lift across green, yellow, orange and red is reproducible at pair-instant level, but it is weaker per encounter: the red lift falls from 3.11 per pair-instant to 2.03 with the first observation of each encounter, about a third lower, and this first-observation value is the signal available when an encounter first appears (Table~\ref{tab:ladder}). The decomposition and the analysis by position (Table~\ref{tab:position}) indicate that the main source is the repeated observation of the same encounter: later observations have higher label rates, and counting each as an independent example is analogous to temporal pseudoreplication. Consequently, a colour with many repeated observations per encounter may appear better supported even when its independent encounter count is modest. Clearance multiplicity is secondary: the median of 65 pair-instants per linked clearance reflects the 350 NM candidate radius (it is 2 among non-green pairs), while 27--41\% of the clearances linked to non-green episodes are linked to more than one episode, depending on the radius. This is why the labels are examined per encounter, keeping all many-to-many links rather than selecting a ``true'' motivating pair without evidence.

A possible estimand is the probability that at least one eligible clearance is associated with an episode within a prespecified horizon, conditional on its entry geometry and available traffic context. A different, aircraft-level estimand would be the probability of any clearance to the focal aircraft. They answer different research questions and require different exposure denominators. We do not argue for one universally correct choice; the aim is to show how the choice affects interpretation.

\subsection{How the four levels should be interpreted}
The calibrated four-colour partition is a convenient summary of the original pair-instant analysis, with increasing positive-label rates, and its thresholds are specific to this dataset and pipeline. Placing $1\le\mmin<1.2$ in green makes the scheme non-monotone near the threshold: an episode near $m=1.1$ can be green while a more widely separated one is yellow, and the present data do not distinguish the possible reasons. The number of levels was chosen through week-based predictive comparison and a simplicity preference, not discovered as an intrinsic property of the airspace. The contribution therefore does not rest on the four levels but on making explicit the data decisions behind such a classification: the unit of evaluation, clearance multiplicity and surveillance coverage. The comparison with continuous and monotone reference models is given in Appendix~\ref{app:models}.

\subsection{Censoring, selection and non-causal interpretation}\label{sec:limits}
Retrospective evaluability is governed mainly by the extent of the recorded surveillance domain. Table~\ref{tab:evidence} shows why 4,307 of the 6,051 red episodes (71.2\%) cannot be followed for the full 20-min horizon. In almost all of them the track of at least one aircraft ends inside FL290--FL410 before the endpoint, and no unevaluable episode is explained by a gap that later closes, so none of these cases is a matching failure. To test whether such track ends coincide with the limit of the recorded coverage, we built an empirical footprint of the data (occupancy grid of about 11 km cells with at least 100 position reports, using all positions of the aircraft in the non-green episodes) and measured the distance from the last position of each track end to the nearest empty cell. The median is 6 NM, the grid resolution, whereas positions of aircraft in evaluable episodes at the anchor time lie a median 85 NM from the edge. The tracks thus end where the recorded coverage ends, about nine minutes (median) before the endpoint. The footprint is empirical and not an operational airspace polygon, so this supports but does not prove geographic exit.

\begin{table}[tbp]
\caption{Evidence for the follow-up of the 6,051 red episodes. Median is the distance from the position to the nearest empty cell of the empirical footprint of recorded positions (cells of about 11 km, so 6 NM is the minimum nonzero value); for evaluable episodes it refers to the positions of both aircraft at the anchor time, and for track ends to the last position of the aircraft whose track ends first. The range in parentheses varies the occupancy threshold from 20 to 500 reports per cell.}\label{tab:evidence}
\centering\small\setlength{\tabcolsep}{4pt}
\begin{tabular}{@{}p{0.36\linewidth}rrrp{0.2\linewidth}@{}}
\toprule
Evidence category & Episodes & \% of red & Median (NM) & Within 12 NM\\
\midrule
Evaluable (observed to the endpoint) & 1,744 & 28.8 & 85 & 6\%\\
Track ends before endpoint, inside FL290--FL410 & 3,823 & 63.2 & 6 & 98.9\% (93.1--99.6\%)\\
Last observed outside FL290--FL410 & 479 & 7.9 & -- & --\\
No observed start state & 5 & 0.1 & -- & --\\
\bottomrule
\end{tabular}
\end{table}

Evaluability also depends on class (28.8\% for red, 64.2\% for orange and 50.6\% for yellow), so complete-case estimates describe episodes that remain observable for 20 min, not all episodes of the class. The audit therefore covers observable encounters; missingness analyses such as stratified models, inverse-probability weighting \cite{seaman2013} with diagnostics, and worst/best-case bounds are natural extensions.

\subsection{Relation to operational MTCD and controller decision support}
EUROCONTROL's MTCD specification already encompasses aircraft encounters, probabilities and interactions \cite{eurocontrol2017}. The four-level indicator is not proposed as a replacement for those functions. The reconstructed associations can be used to examine whether a compact activity indicator provides supplementary information in a future ATM decision-support model; any statement on alert burden would require de-duplicated events per sector and hour rather than raw pair-instants.

\subsection{Limitations}\label{sec:limitations}
The study is retrospective and uses one dataset, 13 sampled weeks of Swedish upper en-route traffic; the levels and their thresholds are specific to this data and pipeline and have not been tested elsewhere. Clearances are linked to encounters by time and aircraft membership, so the labels measure clearance activity, not controller motivation, risk perception or prevented losses of separation. Three data choices bound the interpretation: the label covers eligible heading and level clearances only (a broader definition changes the class ordering, Section~\ref{sec:results_sensitivity}); SCAT excludes flights in incident reports \cite{nilsson2023}, so the most consequential events may be absent; and follow-up is limited by the recorded surveillance domain (Section~\ref{sec:limits}). The forecast is a simplified route-informed predictor, the geometric thresholds are study parameters rather than sector-specific minima, and traffic is stratified by global count only. The test weeks were used repeatedly during development, and the model comparison of Appendix~\ref{app:models} uses episode features, in particular duration, that are known only when the episode ends, so it does not describe performance available at the start of an encounter. The near-threshold green band ($1\le\mmin<1.2$), in particular for $\tau<2$ min where exploratory cells showed elevated rates, was not examined with a dedicated sensitivity analysis and is left for future work.

Three directions follow from these results. First, repeated observation of the same encounter could be handled explicitly with discrete-time hazard models, in which each encounter contributes until a clearance occurs or the encounter ends, instead of counting every observation as an independent example. Second, labels could be built at decision time, using only the information available at the first observation of each encounter, or from the clearance's side, by recording the encounters present when each clearance was issued. Third, the practical consequence of the unit could be measured directly by training the same conflict-resolution or detection model with observation-level and encounter-level labels and comparing their encounter-level performance. Extending the analysis to other datasets, airspace types and geometric indicators would show whether the concentration of the effect in predicted losses of separation is general, and missingness analyses would address the limit set by surveillance coverage.

\section{Conclusions}\label{sec:conclusion}
Counting every observation of an encounter, instead of the encounter, overstates how strongly predicted geometry is associated with a later controller clearance. For the red class, the lift over green falls from 3.11 per pair-instant to 2.03 when only the first observation of each encounter is used, about a third lower (Table~\ref{tab:ladder}). The reduction concentrates in predicted losses of separation, appears with or without the calibrated levels, and comes mainly from the repeated observation of the same encounter, which is analogous to temporal pseudoreplication (Table~\ref{tab:position}).

Only 28.8\% of red episodes can be followed to an observed outcome, and 27--41\% of the clearances linked to non-green episodes relate to more than one episode, so labels derived from clearances cannot be read as independent, uniquely attributed evidence.

We propose that clearance-based labels be reported together with three checks: unit of evaluation, clearance multiplicity and surveillance coverage (Table~\ref{tab:checks}). The findings come from one dataset and are descriptive; the directions that follow from them are listed in Section~\ref{sec:limitations}.


%=================================================================
% MDPI back matter: confirm all declarations with the three coauthors before submission.
\authorcontributions{Conceptualization, B.R., F.H. and E.G.; methodology, B.R., F.H. and E.G.; software, B.R.; validation, B.R.; formal analysis, B.R.; investigation, B.R.; data curation, B.R.; writing---original draft preparation, B.R.; writing---review and editing, B.R., F.H. and E.G.; visualization, B.R.; supervision, F.H. and E.G. \tbc{Confirm that all authors have read and approved the final version of the manuscript.}}

\funding{This research received no external funding.}

\institutionalreview{Not applicable.}
\informedconsent{Not applicable.}

\dataavailability{The SCAT dataset analysed here is publicly available on Mendeley Data (\url{https://doi.org/10.17632/8yn985bwz5.1}) \cite{nilsson2023}. The derived pair-instant and episode tables are stored as Apache Parquet files and amount to several gigabytes; they are not redistributed, and can be regenerated from the SCAT data by running the scripts of the code repository (release v1.1.0). Analysis code and figure-generation scripts are available at \url{https://github.com/BorjaFrade/ai-based-air-traffic-management} (release v1.1.0: \url{https://github.com/BorjaFrade/ai-based-air-traffic-management/releases/tag/v1.1.0}).}

\acknowledgments{During the preparation of this manuscript and the analysis scripts, the authors used generative AI assistance (Claude Sonnet 5.5, Anthropic) to improve the readability of the manuscript and to generate some of the Python scripts used for the data analysis. It was not used to produce the data or the numerical results. The authors reviewed and verified all analyses, numbers, figures and references and take full responsibility for the content. \tbc{Confirm the wording with the coauthors and MDPI policy.}}

\conflictsofinterest{The authors declare no conflicts of interest.}


\abbreviations{Abbreviations}{%
The following abbreviations are used in this manuscript:\\


\noindent
\begin{tabular}{@{}ll}
ADS-B    & Automatic Dependent Surveillance--Broadcast\\
ATC      & Air Traffic Control\\
ATM      & Air Traffic Management\\
FL       & Flight Level\\
ICAO     & International Civil Aviation Organization\\
MTCD     & Medium-Term Conflict Detection\\
NM       & Nautical Mile\\
ROC      & Receiver Operating Characteristic\\
SCAT     & Swedish Civil Air Traffic Control\\
UTC      & Coordinated Universal Time\\
\end{tabular}

}

%%%%%%%%%%%%%%%%%%%%%%%%%%%%%%%%%%%%%%%%%%
\appendixtitles{yes}
\appendixstart
\appendix
\section{Model comparison at the episode level}\label{app:models}
This appendix compares the four-level partition (M0) with reference models on the episode cohort with green observations restricted to $\mmin<2$ (Table~\ref{tab:cohorts}). Three analysis choices (the green inclusion limit $\mmin<2$, the deferral of the Gaussian conflict-probability model M6, and the one-standard-error decision rule) were fixed before the first run of this comparison.

\emph{Splits.} All fitting and model selection use the nine calibration weeks (with leave-one-week-out folds for hyperparameters). Thresholds, parameters and model selection use calibration weeks only; the test weeks are used to score the models that were fitted on them.

\subsection{Models and metrics}
Table~\ref{tab:models} defines M0--M5. M0 converts each frozen colour to its calibration-week empirical clearance rate. M1 similarly converts the binary $\mmin<1$ screen to calibration-week probabilities. M2 and M3 fit logistic probabilities from $\mmin$ alone or cubic spline terms of $(\mmin,\tau)$. M4 adds tertiles of the number of episode starts observed during the preceding hour, prior-clearance status and log episode duration. M5 uses geometry-ordered, isotonic-pooled $\mmin$ quantile partitions with $K=3,4,5$. All models share the same cohort and outcome; spline knots, quantiles, scaling and probability estimates are fitted only on calibration weeks. Logistic regularisation is selected by leave-one-week-out calibration Brier score.

\begin{table}[t]
\caption{Models compared at the episode level. All are fitted on the nine calibration weeks and frozen before the four evaluation weeks are scored.}\label{tab:models}
\begin{tabular}{@{}lp{0.45\linewidth}p{0.30\linewidth}@{}}
\toprule
ID & Model & Role\\
\midrule
M0 & Frozen four-level partition (Eq.~\eqref{eq:color}) & Object of the audit\\
M1 & Binary geometric screen $\mmin<1$ within 20 min & Nominal-detection baseline\\
M2 & Logistic regression on $\mmin$ only & Parsimonious continuous baseline\\
M3 & Logistic regression on smooth terms of $(\mmin,\tau)$ & Continuous two-variable model\\
M4 & M3 plus hourly density stratum, prior clearance and log duration & Contextual model\\
M5 & Monotone-constrained partitions with $K=3,4,5$ levels & Discrete alternatives\\
\bottomrule
\end{tabular}
\end{table}

Models are compared with proper scoring rules that evaluate the predicted probability $\hat p_e$ of the clearance outcome $Y_e\in\{0,1\}$ of each held-out episode $e$ \cite{gneiting2007}. A scoring rule is \emph{proper} when its expected value is optimised by reporting the true probability, so it rewards honest and well-calibrated probabilities rather than merely correct rankings. With $N$ test episodes and weights $w_e$ (all equal to one for one-episode-one-vote, and the number of pair-instants of the episode for pair weighting), the Brier score is the weighted mean squared error of the probabilities and the logarithmic score is the weighted mean negative log-likelihood:
\begin{equation}\label{eq:brier}
 \mathrm{Brier}=\frac{\sum_e w_e\,(\hat p_e-Y_e)^2}{\sum_e w_e},\qquad
 \mathrm{LogScore}=-\frac{\sum_e w_e\left[Y_e\log\hat p_e+(1-Y_e)\log(1-\hat p_e)\right]}{\sum_e w_e}.
\end{equation}
Both are lower for better forecasts: a Brier score of 0 means perfect probabilities, a constant forecast equal to the base rate $\bar y$ scores $\bar y(1-\bar y)$, and the logarithmic score penalises confident wrong probabilities more heavily (probabilities are clipped to $[10^{-6},1-10^{-6}]$). The Brier score is the main metric because it is bounded and its differences are easy to interpret as changes in mean squared probability error. The quantity reported in the comparison tables is $\Delta\mathrm{Brier}=\mathrm{Brier}_{\text{model}}-\mathrm{Brier}_{\mathrm{M0}}$, so negative values favour the alternative model over the frozen colours.

Calibration is summarised by the intercept $a$ and slope $b$ of a logistic regression of the outcome on the model's own log-odds,
\begin{equation}\label{eq:calib}
 \ln\!\frac{\mathbb{P}(Y_e=1)}{1-\mathbb{P}(Y_e=1)}=a+b\,\ln\!\frac{\hat p_e}{1-\hat p_e},
\end{equation}
where $\mathbb{P}(Y_e=1)$ is the probability that episode $e$ is followed by an eligible clearance, fitted on the test episodes. Perfect calibration gives $a=0$ and $b=1$; $b>1$ means the probabilities are too close together (under-confident) and $b<1$ that they are too extreme (over-confident). Decile reliability curves compare mean predicted probability and observed clearance rate in ten bins of $\hat p_e$. The area under the receiver operating characteristic (ROC) curve, the probability that a randomly chosen episode with a clearance is ranked above one without, is secondary because it ignores calibration. Paired Brier uncertainty uses 400 day-cluster bootstrap replicates \cite{field2007}: whole days are resampled with replacement and the paired difference $\Delta\mathrm{Brier}$ is recomputed each time, so that episodes of the same day, which share traffic and controller conditions, are not treated as independent. Four week-level paired differences are descriptive, not significance tests. Alert burden is counted as deduplicated non-green episodes per observed hour, not as raw pair-instants. A Gaussian conflict-probability model is reserved for future work because it targets conflict rather than the recorded clearance outcome.

\emph{Decision rule.} The best of M2--M4 is selected by leave-one-week-out calibration Brier score. It becomes the primary model if its test Brier improvement over M0 exceeds one day-bootstrap standard error of the paired difference. Otherwise M0 remains the recommended summary. The test weeks do not select model complexity.



\subsection{Cohort and window rates}
The common cohort contains 67,367 episodes, including 21,281 in the four test weeks. With the fixed start-based window, test positive rates rise from 16.2\% in green to 21.5\%, 27.7\% and 30.9\% in yellow, orange and red. Rates under the any-observation exposure window are higher in every class (Table~\ref{tab:windows}); this shows how the association window changes the measured target. The colour ordering is preserved under both definitions.


\begin{table}[t]
\caption{Test-week eligible-clearance rates per episode under fixed 300-s exposure from episode start and under a 300-s window after any episode observation.}\label{tab:windows}
\begin{tabular}{@{}lrrr@{}}\toprule
Colour & Episodes & Fixed start window & Any-observation window\\\midrule
Green & 3,943 & 16.2\% & 26.0\%\\
Yellow & 12,548 & 21.5\% & 30.5\%\\
Orange & 2,927 & 27.7\% & 41.0\%\\
Red & 1,863 & 30.9\% & 45.9\%\\\bottomrule
\end{tabular}
\end{table}


\subsection{Results}
M4 was selected among M2--M4 using leave-one-week-out Brier score on the nine calibration weeks. On test it improved the episode-weighted Brier score from 0.1712 for M0 to 0.1690; the paired difference was $-0.00220$ (95\% day-bootstrap interval $[-0.00298,-0.00149]$, bootstrap standard error 0.00040). This exceeds the one-standard-error threshold of the decision rule. The pair-instant-weighted comparison also favoured M4 (0.1482 versus 0.1467; difference $-0.00151$). The paired episode-weighted improvement is negative in each of the four test weeks ($-0.00136,-0.00347,-0.00066,-0.00313$ in chronological order); four weeks are too few for a significance claim. The 17,338 test non-green episodes occurred in 605 hours containing at least one such episode, averaging 28.7 deduplicated episodes per active hour. Most of this gain is already present in M3, which uses only geometry and lead time ($\Delta$Brier $-0.00162$ relative to M0, Table~\ref{tab:bench}). Refitting M3 with each extra term of M4 shows that log episode duration accounts for most of the remaining difference: adding it alone changes the Brier score by $-0.00051$ relative to M3 (95\% day-bootstrap interval $[-0.00067,-0.00034]$), against $-0.00058$ for the full M4. Adding the hourly density or the prior-clearance indicator alone changes it by $+0.00004$ and $-0.00001$, with intervals that include zero, and both together by $+0.00003$. These are comparisons between refitted models, not additive individual effects. Because duration is known only when the episode ends, this part of the improvement is retrospective and is not performance available at $T_0$. Conversely, the four-level partition is not uniformly worse than the continuous alternatives: it scores better than the binary screen M1 and than the logistic model on $\mmin$ alone (M2), and equals the five-level partition (M5, $K=5$). These gains concern the prediction of a subsequent eligible clearance, not conflict detection or safety benefit.

\begin{table}[t]
\caption{Held-out episode-level comparison on the four evaluation weeks (one episode, one vote). $\Delta$Brier is relative to M0; its 95\% day-cluster interval is plotted in Fig.~\ref{fig:benchmarks}.}\label{tab:bench}
\begin{tabular}{@{}lcccc@{}}
\toprule
Model & Brier & Log score & Calib. slope & $\Delta$Brier vs M0\\
\midrule
M0 & 0.1712 & 0.5251 & 1.063 & --\\
M1 & 0.1717 & 0.5265 & 1.156 & $+0.00044$\\
M2 & 0.1722 & 0.5278 & 1.206 & $+0.00092$\\
M3 & 0.1696 & 0.5206 & 0.899 & $-0.00162$\\
M4 & 0.1690 & 0.5189 & 0.934 & $-0.00220$\\
M5 ($K=3$) & 0.1717 & 0.5265 & 1.156 & $+0.00044$\\
M5 ($K=4$) & 0.1717 & 0.5265 & 1.151 & $+0.00043$\\
M5 ($K=5$) & 0.1712 & 0.5252 & 1.130 & $-0.00002$\\
\bottomrule
\end{tabular}
\end{table}

\begin{figure}[tbp]
\centering
\subfloat[Held-out episode comparison: paired Brier differences versus M0.]{%
\begin{minipage}[t]{0.49\linewidth}
\centering
\subimg{0.42}{figures/fig08a_brier.png}
\end{minipage}%
}\hfill
\subfloat[Reliability on test episodes.]{%
\begin{minipage}[t]{0.49\linewidth}
\centering
\subimg{0.42}{figures/fig08b_reliability.png}
\end{minipage}%
}
\caption{Held-out episode-weighted comparison of M0--M5. Negative differences favour the candidate over M0. Error bars are 95\% day-cluster bootstrap intervals from 400 replicates. Reliability bins are deciles of predicted probability (duplicates collapsed); only M0 and M2--M4 are drawn in (b).}\label{fig:benchmarks}
\end{figure}


\subsection{Sensitivity}
With a fixed 120-s later window, M4 remains the calibration-selected continuous model and improves Brier over M0 by 0.00161. At 600 s, class-rate ordering is restored and M4 improves Brier by 0.00220. Reducing the prior-clearance window from 600 to 300 s leaves M4 preferred, with a 0.00213 Brier improvement. With 240- and 480-s episode gaps, the class-rate ordering is retained and M4 remains selected, improving Brier by 0.00156 and 0.00251, respectively. With the green inclusion limit relaxed to $\mmin<3$, M4 remains calibration-selected and improves Brier by 0.00138. With no green limit, fitting the primary M0 and M4 specifications on the enlarged calibration cohort yields test Brier scores of 0.13355 and 0.13247, respectively; this large-cohort sensitivity freezes M4's primary regularisation choice rather than repeating the full model-selection grid. Relabelling the same episodes with all 908,165 extracted heading/level records raises clearance rates sharply: in test they are 23.6\%, 36.7\%, 58.3\% and 34.5\% for green through red. Red now falls below orange. M4 remains calibration-selected and improves Brier by 0.04181 under this substantially different target; the magnitude is not comparable with the primary experiment, because the target itself changes.

\section{Observed manoeuvres and inertial screening}\label{app:manoeuvres}
The screening also checks whether an observable heading or level change is compatible with the recorded clearance and records its approximate onset. Finally, it compares observed separation with a constant-horizontal-and-vertical-velocity projection based on pre-clearance states. This projection is a sensitivity screen only: it ignores planned manoeuvres and later actions of either aircraft, so it cannot estimate prevented losses of separation.

Among 28,324 distinct clearances linked to non-green encounters, 16,566 (58.5\%) showed a trajectory change compatible with the recorded heading or level target, while 23,875 (84.3\%) showed some observable change, whether or not it could be matched confidently to the clearance. A further 3,071 had no observable execution and 1,378 were indeterminate. For compatible changes, the proxy onset quantiles were 20, 30 and 110 seconds after issue. These are compatibility categories, not a reconstruction of controller intent. A large fraction of linked clearances (22,159) had another recorded clearance within five minutes.

An exploratory constant-velocity pre-clearance projection had sufficient pre-clearance state and observed follow-up for 30,081 of 63,524 projection rows. Among the rows compatible with a manoeuvre, 578 had projected $m<1$ and fully observed $m\ge1$; only 17 exceeded a 90th-percentile absolute projection-error threshold computed from 3,807 full-coverage calibration reference rows without linked clearances. Nearby clearances or interactions may explain even those 17 cases. In test, all three cases above that error threshold had a nearby clearance. The screening therefore does not provide a count of \emph{conflicts prevented by clearance}.




%%%%%%%%%%%%%%%%%%%%%%%%%%%%%%%%%%%%%%%%%%
\begin{adjustwidth}{-\extralength}{0cm}
%\MDPIprintendnotes % Un-comment to print endnotes (heading: Note / Notes by count)

\reftitle{References}

% References must be numbered in order of appearance in the text (including citations in tables and legends) and listed individually at the end of the manuscript. We recommend preparing the references with a bibliography software package, such as EndNote, ReferenceManager or Zotero, to avoid typing mistakes and duplicated references. Include the digital object identifier (DOI) for all references where available.
% Please provide either the correct journal abbreviation (e.g. according to the “The list of Title Word Abbreviations” https://portal.issn.org/ltwa) or the full name of the journal. 
% Citations and references in the Supplementary Materials are permitted provided that they also appear in the reference list here. 
% In the text, reference numbers should be placed in square brackets [ ] and placed before the punctuation; for example, [1], [1–3] or [1,3]. For embedded citations in the text with pagination, use both parentheses and brackets to indicate the reference number and page numbers; for example, [5] (p. 10), or [6] (pp. 101–105).

%=====================================
% References, variant A: external bibliography
%=====================================
% \bibliography{your_external_BibTeX_file}

%=====================================
% References, variant B: internal bibliography
%=====================================

% ACS format
\begin{thebibliography}{999}

% Reference 1 - Technical report
\bibitem{eurocontrol2017}
EUROCONTROL.
\textit{EUROCONTROL Specification for Medium-Term Conflict Detection (MTCD)},
2nd ed.; EUROCONTROL-SPEC-0139; EUROCONTROL, 2017.
Available online:
\url{https://www.eurocontrol.int/publication/eurocontrol-specification-medium-term-conflict-detection-mtcd}.

% Reference 2 - Journal article
\bibitem{kuchar2000}
Kuchar, J. K.; Yang, L. C.
A Review of Conflict Detection and Resolution Modeling Methods.
{\em IEEE Trans. Intell. Transp. Syst.}
{\bf 2000}, {\em 1} (4), 179--189.
\url{https://doi.org/10.1109/6979.898217}.

% Reference 3 - Journal article
\bibitem{paielli1997}
Paielli, R. A.; Erzberger, H.
Conflict Probability Estimation for Free Flight.
{\em J. Guid. Control Dyn.}
{\bf 1997}, {\em 20} (3), 588--596.
\url{https://doi.org/10.2514/2.4081}.

% Reference 4 - Journal article
\bibitem{prandini2000}
Prandini, M.; Hu, J.; Lygeros, J.; Sastry, S.
A Probabilistic Approach to Aircraft Conflict Detection.
{\em IEEE Trans. Intell. Transp. Syst.}
{\bf 2000}, {\em 1} (4), 199--220.
\url{https://doi.org/10.1109/6979.898224}.

% Reference 5 - Book / Technical publication
\bibitem{icao4444}
International Civil Aviation Organization.
\textit{Procedures for Air Navigation Services --
Air Traffic Management (PANS-ATM)},
16th ed.; Doc 4444;
ICAO: Montr\'eal, Canada, 2016.
Available online:
\url{https://s-store.icao.int/en/procedures-for-air-navigation-services-air-traffic-management-doc-4444}.

% Reference 6 - Journal article
\bibitem{bastas2022}
Bastas, A.; Vouros, G. A.
Data-Driven Prediction of Air Traffic Controllers
Reactions to Resolving Conflicts.
{\em Inf. Sci.}
{\bf 2022}, {\em 613}, 763--785.
\url{https://doi.org/10.1016/j.ins.2022.09.015}.

% Reference 7 - Journal article
\bibitem{bastas2023}
Bastas, A.; Vouros, G. A.
Data-Driven Modeling of Air Traffic Controllers'
Policy to Resolve Conflicts.
{\em Aerospace}
{\bf 2023}, {\em 10} (6), 557.
\url{https://doi.org/10.3390/aerospace10060557}.

% Reference 8 - Journal article
\bibitem{delgado2022}
Delgado-Aguilera Jurado, R.;
G\'omez Comendador, V. F.;
Zamarre\~no Su\'arez, M.;
P\'erez Moreno, F.;
Verdonk Gallego, C. E.;
Arnaldo Vald\'es, R. M.
Assessment of Potential Conflict Detection by the ATCo.
{\em Aerospace}
{\bf 2022}, {\em 9} (9), 522.
\url{https://doi.org/10.3390/aerospace9090522}.

% Reference 9 - Journal article
\bibitem{nilsson2023}
Nilsson, J.; Unger, J.
Swedish Civil Air Traffic Control Dataset.
{\em Data Brief}
{\bf 2023}, {\em 48}, 109240.
\url{https://doi.org/10.1016/j.dib.2023.109240}.

% Reference 10 - Conference proceeding
\bibitem{renkhoff2025}
Renkhoff, J.
Identifying Potential Conflicts and Corresponding
Resolutions in the Upper Airspace from Historical
Air Traffic Data.
In Proceedings of the
{\em Deutscher Luft- und Raumfahrtkongress (DLRK) 2025};
2025.
\url{https://doi.org/10.25967/650174}.

% Reference 11 - Journal article
\bibitem{renkhoff2026}
Renkhoff, J.; W\"ustenbecker, N.;
Jameel, M.; Schier-Morgenthal, S.
Predicting and Reconstructing Clearances from
Air Traffic Data Using a Supervised Learning Approach.
{\em CEAS Aeronaut. J.}
{\bf 2026}, {\em 17}, 343--354.
\url{https://doi.org/10.1007/s13272-025-00854-x}.

% Reference 12 - Journal article
\bibitem{wustenbecker2025}
W\"ustenbecker, N.; Renkhoff, J.;
Zeppenfeld, D.; Jameel, M.;
Schier-Morgenthal, S.
Analysis and Prediction of Pilot Response Time
to Air Traffic Control Clearances.
{\em CEAS Aeronaut. J.}
{\bf 2026}, {\em 17}, 507--520.
\url{https://doi.org/10.1007/s13272-025-00848-9}.

% Reference 13 - Journal article
\bibitem{guleria2024}
Guleria, Y.; Pham, D.-T.; Alam, S.;
Tran, P. N.; Durand, N.
Towards Conformal Automation in Air Traffic Control:
Learning Conflict Resolution Strategies through
Behavior Cloning.
{\em Adv. Eng. Inform.}
{\bf 2024}, {\em 59}, 102273.
\url{https://doi.org/10.1016/j.aei.2023.102273}.

% Reference 14 - Journal article
\bibitem{gaume2025}
Gaume, K.; Olive, X.; Gianazza, D.;
Alligier, R.; Durand, N.
Extracting Aircraft Conflict-Resolution Situations
from Historical ADS-B Data.
{\em Transp. Res. Interdiscip. Perspect.}
{\bf 2025}, {\em 34}, 101669.
\url{https://doi.org/10.1016/j.trip.2025.101669}.

% Reference 15 - Journal article
\bibitem{gaume2023}
Gaume, K.; Olive, X.; Gianazza, D.;
Alligier, R.; Durand, N.
A Catalogue of Deconfliction Actions Extracted
from Historical ADS-B Data.
{\em J. Open Aviat. Sci.}
{\bf 2023}, {\em 1} (2).
\url{https://doi.org/10.59490/joas.2023.7222}.

% Reference 16 - Journal article
\bibitem{xia2024}
Xia, C.; Verdonk Gallego, C. E.;
Bracero, A. F.; G\'omez Comendador, V. F.;
Arnaldo Vald\'es, R. M.
Trajectory Predictor and Conflict Detection Figures
of Merit for a Performance-Based Adaptive
Air Traffic Monitoring System.
{\em Aerospace}
{\bf 2024}, {\em 11} (2), 155.
\url{https://doi.org/10.3390/aerospace11020155}.

% Reference 17 - Journal article
\bibitem{hurlbert1984}
Hurlbert, S. H.
Pseudoreplication and the Design of Ecological Field Experiments.
{\em Ecol. Monogr.}
{\bf 1984}, {\em 54} (2), 187--211.
\url{https://doi.org/10.2307/1942661}.

% Reference 18 - Journal article
\bibitem{killick2012}
Killick, R.; Fearnhead, P.; Eckley, I. A.
Optimal Detection of Changepoints with a Linear
Computational Cost.
{\em J. Am. Stat. Assoc.}
{\bf 2012}, {\em 107} (500), 1590--1598.
\url{https://doi.org/10.1080/01621459.2012.737745}.

% Reference 18 - Journal article
\bibitem{field2007}
Field, C. A.; Welsh, A. H.
Bootstrapping Clustered Data.
{\em J. R. Stat. Soc. Ser. B}
{\bf 2007}, {\em 69} (3), 369--390.
\url{https://doi.org/10.1111/j.1467-9868.2007.00593.x}.

% Reference 19 - Journal article
\bibitem{gneiting2007}
Gneiting, T.; Raftery, A. E.
Strictly Proper Scoring Rules, Prediction,
and Estimation.
{\em J. Am. Stat. Assoc.}
{\bf 2007}, {\em 102} (477), 359--378.
\url{https://doi.org/10.1198/016214506000001437}.

% Reference 20 - Journal article
\bibitem{seaman2013}
Seaman, S. R.; White, I. R.
Review of Inverse Probability Weighting
for Dealing with Missing Data.
{\em Stat. Methods Med. Res.}
{\bf 2013}, {\em 22} (3), 278--295.
\url{https://doi.org/10.1177/0962280210395740}.

\end{thebibliography}


% If authors have biography, please use the format below
\section*{Short Biography of Authors}
\bio[bio/borja.jpg]{Borja Rodríguez Frade}{received his degree in Electrical Engineering from the University of A Coruña, Spain, in 2010, and his master's degree in Renewable Energy from the European University, Spain, in 2012. He received his Ph.D. from the University of A Coruña in 2021, with research focused on multiphysics simulation, co-simulation, and model-based testing of automotive electric powertrains. He currently works as a Technical Lead at Indra (Minsait), focusing on numerical modelling, simulation, and software development for electrical power systems. He is also pursuing a master's degree in Artificial Intelligence at UNED, where his current research explores applications of AI to air traffic management.}
\bio[bio/felix.png]{Félix Hernández-del-Olmo}{received his B.Sc. degree in Physics from the Universidad Autónoma de Madrid, Spain, in 1996, and his Engineering degree and Ph.D. in Computer Science from the Universidad Nacional de Educación a Distancia (UNED), Spain, in 1999 and 2005, respectively. He is currently an Associate Professor in the Department of Artificial Intelligence at UNED. His academic work includes machine learning and intelligent systems.}
\bio[bio/elena.png]{Elena Gaudioso}{received her B.Sc. degree in Mathematics from the Universidad Complutense de Madrid, Spain, in 1996, and her Ph.D. in Sciences from the Universidad Nacional de Educación a Distancia (UNED), Spain, in 2003. She is currently an Associate Professor in the Department of Artificial Intelligence at UNED and coordinates its Master's programme in Artificial Intelligence Research. Her research interests include machine learning, recommender systems and adaptive systems.}
%\bio{Firstname Lastname}{Biography of second author without a photo.}

% For the MDPI journals use author-date citation, please follow the formatting guidelines on http://www.mdpi.com/authors/references
% To cite two works by the same author: \citeauthor{ref-journal-1a} (\citeyear{ref-journal-1a}, \citeyear{ref-journal-1b}). This produces: Whittaker (1967, 1975)
% To cite two works by the same author with specific pages: \citeauthor{ref-journal-3a} (\citeyear{ref-journal-3a}, p. 328; \citeyear{ref-journal-3b}, p.475). This produces: Wong (1999, p. 328; 2000, p. 475)

%%%%%%%%%%%%%%%%%%%%%%%%%%%%%%%%%%%%%%%%%%
%% for journal Sci
%\reviewreports{\\
%Reviewer 1 comments and authors’ response\\
%Reviewer 2 comments and authors’ response\\
%Reviewer 3 comments and authors’ response
%}
%%%%%%%%%%%%%%%%%%%%%%%%%%%%%%%%%%%%%%%%%%
\PublishersNote{}
\end{adjustwidth}
\end{document}